\section{Nonlinear Estimates for One Ring}
\label{sec:single}

In this section we study, in ring units, one ring of the approximate solutions: a
ring close to $\Phi$, whose deviation from $\Phi$ is driven by its own
nonlinearity and by the field $h$ of the vorticity inside and outside it, apart
from the point-vortex part of the inner vorticity, which enters the linearization
$\Lin(G)$ (Lemma~\ref{lem:reduction}). We prove three results. A ring that the
instability has not yet reached stays close to $\Phi$, and its deviation equals
the linear response to $h$ up to an error of second order
(Theorem~\ref{thm:quiet}). A ring whose component along
the unstable eigenfunction grows follows that eigenfunction until the envelope of
this component reaches a small size (Theorem~\ref{thm:active}). And the innermost ring of
the approximate datum can be given a small component along the unstable
eigenfunction by a rearrangement of $\Phi$ (Lemma~\ref{lem:seed}).
Sections~\ref{sec:relay} and~\ref{sec:bootstrap} apply Theorems~\ref{thm:quiet}
and~\ref{thm:active} to every ring, in its own units, and Section~\ref{sec:kick} also uses
Theorem~\ref{thm:quiet}. Lemma~\ref{lem:seed} gives the innermost ring of the
datum of Definition~\ref{def:approx}.

Throughout this section we assume \textup{(N1)}--\textup{(N3)}, we work in ring
units, and we use the objects of Definition~I.\ref{I-def:eigendata}: the unstable
eigenvalue $\lamst(G)$ of $\Lin(G)$ in the mode $3$ and its eigenfunction
$\phist(G)$, normalized in $L^2(\R^2)$; the functional $\ellst(G)$ that gives the
coefficient of a function along $\phist(G)$; the Riesz projection $\Riesz(G)$
onto the span of $\phist(G)$ and $\overline{\phist(G)}$; and the growth rates
$\sigma(G)=\Re\lamst(G)$ and $\sigb=\min_{G\in[0,\Gone]}\sigma(G)$. When $G$ is fixed
we write $\lamst=\lamst(G)$, $\phist=\phist(G)$, $\ellst=\ellst(G)$,
$\Riesz=\Riesz(G)$. The proofs of
Theorems~\ref{thm:quiet} and~\ref{thm:active} follow Grenier \cite{Grenier}: an
approximate solution of high order is built from the linear semigroup, and the order is so large that a crude energy estimate for the
remainder suffices.

We use the norms of $H^k$, $W^{1,\infty}$ and $C^k(N')$ fixed in
Section~I.\ref{I-sec:prelim}, and the embedding $\|f\|_{L^\infty}\le\|f\|_{H^2}$ of
Lemma~\ref{lem:embed}\textup{(E2)}. A vector field $h$ on a neighborhood of $N'$
is \textit{$3$-fold rotationally symmetric} if $h(\Orot x)=\Orot h(x)$.

The equation for the deviation of a ring from $\Phi$
(Lemma~\ref{lem:reduction}\textup{(3)} below) assumes that the ring stays in
the annulus $N$, and the continuity arguments of this section keep it there by
bounding how far its material moves radially. Let $\flow$ be the flow
of a solution of Theorem~\ref{thm:smooth}, let $\omega_{\mathrm r}(0)$ have compact
support, let $\omega_{\mathrm r}(t)=\omega_{\mathrm r}(0)\circ\flow_t^{-1}$, and let
$\delta_{\mathrm r}(t)$ be the supremum of $\big||\flow_\tau(y)|-|y|\big|$ over
$\tau\in[0,t]$ and $y\in\supp\omega_{\mathrm r}(0)$ ($0$ if the support is empty);
it is a maximum if the support is nonempty, it is continuous and
nondecreasing, and $\delta_{\mathrm r}(0)=0$, because $(\tau,y)\mapsto\flow_\tau(y)$
is continuous. We say that \textit{every trajectory of $\omega_{\mathrm r}$ moves
radially by at most $\delta$ on $[0,t]$} if $\delta_{\mathrm r}(t)\le\delta$. Since
$\supp\omega_{\mathrm r}(\tau)=\flow_\tau(\supp\omega_{\mathrm r}(0))$, the support
of $\omega_{\mathrm r}(\tau)$, $\tau\le t$, then lies within radial distance $\delta$
of $\supp\omega_{\mathrm r}(0)$.

For $G\in[0,\Gone]$ and an integer $k\ge0$, $e^{t\Lin(G)}$ denotes the semigroup
of Lemma~I.\ref{I-lem:semigroup-construct} for the ring in ring units. It preserves
each mode, hence the subspace $X^k$ of the space $Y^k$ of that lemma
consisting of the functions whose mode components vanish outside $3\Z$. It maps
real-valued functions to real-valued functions, since it is the sum of a
Dyson-Phillips series built from the two operators of that lemma, which both do:
the transport group $e^{t\Lin^{\rm tr}}$, which rotates each circle $\{|y|=r\}$ by
the angle $\Omega(r;G)t$, and the coupling operator $\mathcal K$, here
$\mathcal Kw=-(\KBS\ast w)\cdot\nabla\Phi$. The approximate solutions of this
section are Duhamel integrals of this semigroup. The next lemma shows that they are
solutions of the forced linear equation, differentiable in time as functions with
values in $X^{k-1}$ when the data lie in $X^k$.

\begin{lemma}
\label{lem:duhamel}
Let $G\in[0,\Gone]$, $T>0$, and let $k\ge1$ be an integer. Let $f_0\in X^k$ and
$g\in C([0,T];X^k)$, and put
$u(t)=e^{t\Lin(G)}f_0+\int_0^te^{(t-\tau)\Lin(G)}g(\tau)\dd\tau$. Then
$u\in C([0,T];X^k)\cap C^1([0,T];X^{k-1})$, and $\partial_tu=\Lin(G)u+g$ on
$[0,T]$. If $f_0$ and $g$ are real-valued, so is $u$.
\end{lemma}
\begin{proof}
We first prove the continuity of $u$ in $X^k$, and then its differentiability in
$X^{k-1}$, where the generator of the semigroup acts on $X^k$ as $\Lin(G)$.
The semigroup $S(t)=e^{t\Lin(G)}$ is strongly continuous on $X^k$ and, by the
uniform boundedness principle, bounded on $[0,T]$. So $\tau\mapsto S(t-\tau)g(\tau)$
is continuous into $X^k$: the argument given below for the continuity of
$(s,\tau)\mapsto S(s-\tau)H(\tau)$ in $X^{k-1}$ uses only the strong continuity of
$S$ and a bound on $\|S(t)\|$ for $t$ in a compact interval, and both hold in
$X^k$. Hence
$\int_0^tS(t-\tau)g(\tau)\dd\tau=\int_0^tS(\tau)g(t-\tau)\dd\tau$
is continuous in $t$ by dominated convergence, and $u\in C([0,T];X^k)$. Moreover
$X^k\hookrightarrow X^{k-1}$, and $S$ is a $C_0$-semigroup on $X^{k-1}$, since it
is one on $Y^{k-1}$ and $X^{k-1}=Y^{k-1}_{3\Z}$ is a closed invariant subspace,
both by Lemma~I.\ref{I-lem:semigroup-construct}\textup{(2)}. By
Lemma~I.\ref{I-lem:semigroup-construct}\textup{(2)}, the generator $A$ of $S$ on
$X^{k-1}$ (the part in $X^{k-1}$ of the generator on $Y^{k-1}$ of that lemma)
satisfies $X^k\subset D(A)$, $Ay=\Lin(G)y$ and $\|Ay\|_{X^{k-1}}\le C\|y\|_{X^k}$ for
$y\in X^k$; here $\Lin(G)$ maps $X^k$ into $X^{k-1}$ because it preserves modes
(Lemma~I.\ref{I-lem:sector}\textup{(2)}). Below, $\|\cdot\|$ is the norm of $X^{k-1}$ and
the operator norm on $X^{k-1}$.

All integrals below are Riemann integrals of continuous functions with values
in a Banach space; Fubini's theorem holds for them, because it holds after
applying any continuous linear functional. We extend $g$ by $g(t)=g(T)$ for
$T<t\le2T$. Since each $t\mapsto S(t)x$, $x\in X^{k-1}$, is continuous, hence bounded on $[0,2T]$, the
uniform boundedness principle gives $M_T:=\sup_{0\le t\le2T}\|S(t)\|<\infty$.

Let $x\in D(A)$. For $t\ge0$,
$\epsilon^{-1}(S(\epsilon)-I)S(t)x=S(t)\epsilon^{-1}(S(\epsilon)x-x)\to S(t)Ax$, so $S(t)x\in D(A)$,
$AS(t)x=S(t)Ax$, and $S(t)Ax$ is the right derivative of $t\mapsto S(t)x$. For
$0<\epsilon\le t\le2T$,
$\epsilon^{-1}(S(t)x-S(t-\epsilon)x)-S(t)Ax=S(t-\epsilon)\big(\epsilon^{-1}(S(\epsilon)x-x)-Ax\big)+\big(S(t-\epsilon)-S(t)\big)Ax\to0$,
because $\|S(t-\epsilon)\|\le M_T$. Thus $t\mapsto S(t)x$ lies in $C^1([0,2T];X^{k-1})$ with
derivative $S(t)Ax$, and
\begin{equation}
\label{eq:semigroup-FTC}
S(r)x-x=\int_0^rS(s)Ax\dd s\qquad(x\in D(A),\ 0\le r\le2T).
\end{equation}

Let $H\in C([0,2T];X^{k-1})$. The identity
$S(s'-\tau')H(\tau')-S(s-\tau)H(\tau)=S(s'-\tau')(H(\tau')-H(\tau))
+(S(s'-\tau')-S(s-\tau))H(\tau)$, whose first term has norm at most
$M_T\|H(\tau')-H(\tau)\|$ and whose second tends to $0$ by strong continuity,
shows that $(s,\tau)\mapsto S(s-\tau)H(\tau)$ is continuous on
$\{0\le\tau\le s\le2T\}$. Moreover
$s\mapsto\int_0^sS(s-\tau)H(\tau)\dd\tau$ is continuous on $[0,2T]$: for
$0\le s<s'\le2T$ the difference of its values is
$\int_s^{s'}S(s'-\tau)H(\tau)\dd\tau+\int_0^sS(s-\tau)(S(s'-s)-I)H(\tau)\dd\tau$,
of norm at most $M_T\sup\|H\|(s'-s)+2TM_T\sup_\tau\|(S(s'-s)-I)H(\tau)\|$; and the
operators $S(\epsilon)-I$, $0\le\epsilon\le2T$, are bounded by $M_T+1$ and tend strongly to
$0$ as $\epsilon\to0$, hence uniformly on the compact set $H([0,2T])$.

Now put $V(t)=\int_0^tS(t-\tau)g(\tau)\dd\tau$ and
$W(t)=\int_0^tS(t-\tau)Ag(\tau)\dd\tau$. Since
$\|A(g(\tau)-g(\tau'))\|\le C\|g(\tau)-g(\tau')\|_{X^k}$, $Ag$ is continuous
into $X^{k-1}$, so $W$ is continuous by the previous paragraph with $H=Ag$. By
Equation~\eqref{eq:semigroup-FTC} with $x=g(\tau)\in X^k\subset D(A)$ and Fubini's theorem,
\[
V(t)=\int_0^t\Big(g(\tau)+\int_\tau^tS(s-\tau)Ag(\tau)\dd s\Big)\dd\tau
=\int_0^t\big(g(s)+W(s)\big)\dd s ,
\]
so $V\in C^1([0,2T];X^{k-1})$ and $V'=g+W$. For $t\in[0,T]$ and $0<\epsilon\le T$ we have
$V(t+\epsilon)=S(\epsilon)V(t)+\int_t^{t+\epsilon}S(t+\epsilon-\tau)g(\tau)\dd\tau$, and $\epsilon^{-1}$ times the
last integral tends to $g(t)$ as $\epsilon\to0$, by the continuity of
$(s,\tau)\mapsto S(s-\tau)g(\tau)$ at $(t,t)$. Hence
$\epsilon^{-1}(S(\epsilon)V(t)-V(t))\to V'(t)-g(t)=W(t)$; that is, $V(t)\in D(A)$ and
$AV(t)=W(t)$. With the case $x=f_0\in X^k\subset D(A)$ above, $u=S(\cdot)f_0+V$ lies in
$C^1([0,T];X^{k-1})$, $u(t)\in D(A)$ with $Au(t)=S(t)Af_0+W(t)$, and
$\partial_tu=S(\cdot)Af_0+g+W=Au+g=\Lin(G)u+g$, since $u(t)\in X^k$. Finally, $S(t)$
preserves real-valued functions.
\end{proof}

\subsection{The Perturbation Equation}

We split a smooth solution into a ring $\omega_{\mathrm r}$, the vorticity
$\omega_{\mathrm i}$ inside it and the vorticity $\omega_{\mathrm o}$ outside it,
each transported by the full flow, and write the Euler equations as an equation
for the deviation $w=\omega_{\mathrm r}-\Phi$ of the ring from $\Phi$. The
vorticity inside the ring enters the linearization $\Lin(G)$ as a point vortex at
the origin with its circulation $2\pi G$; the rest of its field, and the field of
the vorticity outside the ring, form the field $h$, which forces the equation.

\begin{lemma}
\label{lem:reduction}
Let $\omega$ be a solution given by Theorem~\ref{thm:smooth} whose datum
$\omega(0)$ is $3$-fold rotationally symmetric, and let $\flow_t$ be its flow. Let
$\omega(0)=\omega_{\mathrm r}(0)+\omega_{\mathrm i}(0)+\omega_{\mathrm o}(0)$ with
pairwise disjoint supports, $\omega_{\mathrm r}(0)$, $\omega_{\mathrm i}(0)$,
$\omega_{\mathrm o}(0)$ each $3$-fold rotationally symmetric, and put
$\omega_\bullet(t)=\omega_\bullet(0)\circ\flow_t^{-1}$ for
$\bullet\in\{\mathrm r,\mathrm i,\mathrm o\}$. Put
$G=\frac1{2\pi}\int\omega_{\mathrm i}(0)$. Let $T>0$ and assume that for every
$t\in[0,T]$, $\supp\omega_{\mathrm i}(t)\subset\overline B_{a/8}$ and
$\supp\omega_{\mathrm o}(t)\cap B_6=\emptyset$. Put
\[
h(t)=\KBS\ast\big(\omega_{\mathrm i}(t)+\omega_{\mathrm o}(t)\big)-G\,\frac{x^\perp}{|x|^2}
\qquad\text{on a neighborhood of }N'.
\]
Then:
\begin{enumerate}[label=\textup{(\arabic*)}]
\item for every $t\in[0,T]$, $h(t)$ is smooth, divergence-free and
  $3$-fold rotationally symmetric on a neighborhood of $N'$; and $t\mapsto h(t)\in C^k(N')$ is
  continuous on $[0,T]$ for every $k$;
\item $\omega_{\mathrm r}\in C^\infty(\R\times\R^2)$, and on $[0,T]\times N'$ the
  velocity $v=\KBS\ast\omega$ is
  \[
  v=\Omega(|x|;G)\,x^\perp+h+\KBS\ast(\omega_{\mathrm r}-\Phi),
  \]
  with $\Omega$ of Definition~I.\ref{I-def:sector-setting} for $(\Phi,G)$; consequently,
  if $t\in[0,T]$ and $\flow_\tau(y)\in N'$ for every $\tau\in[0,t]$, then
  $|\flow_t(y)|-|y|=\int_0^t\big((h+\KBS\ast(\omega_{\mathrm r}-\Phi))\cdot\er\big)(\tau,\flow_\tau(y))\dd\tau$;
\item if $\supp\omega_{\mathrm r}(t)\subset N$ for $t\in[0,T]$, then the function
  $w=\omega_{\mathrm r}-\Phi$ is real-valued, it lies in $C^1([0,T];X^k)$ for
  every $k$, and
  \begin{equation}
  \label{eq:ring}
  \partial_tw=\Lin(G)w+B(w,w)-h\cdot\nabla w-h\cdot\nabla\Phi,\qquad B(v,w)=-(\KBS\ast v)\cdot\nabla w,
  \end{equation}
  on $[0,T]\times\R^2$, where $h\cdot\nabla w$ and $h\cdot\nabla\Phi$ denote the
  functions equal to these products on $N'$ and to $0$ elsewhere.
\end{enumerate}
\end{lemma}
\begin{proof}
Put $K_\bullet=\supp\omega_\bullet(0)$. These compact sets are pairwise
disjoint, so each $\omega_\bullet(0)$ agrees with $\omega(0)$ near $K_\bullet$ and
vanishes near the other two; hence $\omega_\bullet(0)\in C_c^\infty(\R^2)$. The
map $(t,x)\mapsto\flow_t^{-1}(x)$ is the value at time $0$ of the solution of
$\dot z=v(\tau,z)$ with $z(t)=x$. Since $v$ is smooth and bounded on
$[-T',T']\times\R^2$ for every $T'>0$ (Theorem~\ref{thm:smooth}\textup{(2)}),
Lemma~\ref{lem:flow}\textup{(a)}, \textup{(b)} (Appendix~\ref{app:smooth}) shows that this map
is defined and
$C^\infty$ on $\R\times\R^2$. Hence $\omega_\bullet\in C^\infty(\R\times\R^2)$,
$\supp\omega_\bullet(t)=\flow_t(K_\bullet)$, and $\int\omega_{\mathrm i}(t)=2\pi G$
because $\flow_t$ preserves area. By Theorem~\ref{thm:smooth}\textup{(6)},
$\omega(t)$ is $3$-fold rotationally symmetric. Since $\KBS(\Orot x)=\Orot\KBS(x)$, we get
$(\KBS\ast g)(\Orot x)=\Orot(\KBS\ast g)(x)$ for every $3$-fold rotationally symmetric
$g\in (L^1\cap L^\infty)(\R^2)$; so $v(t,\Orot x)=\Orot v(t,x)$, uniqueness for the flow
gives $\flow_t(\Orot y)=\Orot\flow_t(y)$, and each $\omega_\bullet(t)$ is $3$-fold
rotationally symmetric.

\emph{Proof of (1).} By continuity of the flow on compact sets, and since
$\flow_t(K_{\mathrm i})=\supp\omega_{\mathrm i}(t)\subset\overline B_{a/8}$,
$\rho_{\mathrm i}:=\max\{|\flow_t(y)|:t\in[0,T],\,y\in K_{\mathrm i}\}\le a/8<a/4$
($\rho_{\mathrm i}:=0$ if $K_{\mathrm i}=\emptyset$), and $|\flow_t(y)|\ge6$ for
$t\in[0,T]$ and $y\in K_{\mathrm o}$. So the open annulus
$U=\{\frac12(\rho_{\mathrm i}+\frac a4)<|x|<4\}\supset N'$ has distance at least
$\min(\frac12(\frac a4-\rho_{\mathrm i}),2)>0$ from
$\supp(\omega_{\mathrm i}(t)+\omega_{\mathrm o}(t))$ for $t\in[0,T]$. Since
$|\partial^\beta\KBS(z)|\le C_\beta|z|^{-1-|\beta|}$, and $\omega_{\mathrm i}$,
$\omega_{\mathrm o}$ are continuous with supports in a fixed compact set for
$t\in[0,T]$ (Theorem~\ref{thm:smooth}\textup{(1)}), we may differentiate under the
integral sign: $h(t)\in C^\infty(U)$, every $\partial^\beta_xh$ is continuous on
$[0,T]\times U$, and in particular $t\mapsto h(t)\in C^k(N')$ is continuous. The
fields $x\mapsto\KBS(x-y)=\nabla^\perp_x\big(\frac1{2\pi}\log|x-y|\big)$ and
$x^\perp/|x|^2$ are divergence-free away from $y$ and from $0$; hence
so is $h(t)$ on $U$, and it is $3$-fold rotationally symmetric by the equivariance above and
$(\Orot x)^\perp=\Orot x^\perp$.

\emph{Proof of (2).} For a radial $g\in C_c^\infty(\R^2)$, $\KBS\ast g$ is
equivariant under every rotation, so $(\KBS\ast g)\cdot\er$ and
$(\KBS\ast g)\cdot\etheta$ are constant on each $C_r=\{|x|=r\}$, $r>0$. Fix $r>0$,
orient $C_r$ counterclockwise, and let $|y|\ne r$. Along $C_r$,
$\KBS(x-y)\cdot\er\,\dd s=-\frac1{2\pi}\dd\log|x-y|$ and
$\KBS(x-y)\cdot\etheta\,\dd s=\frac1{2\pi}\dd\arg(x-y)$; the first integrates to
$0$, the second to the winding number of $C_r$ around $y$, which is $1$ if
$|y|<r$ and $0$ if $|y|>r$. Integrating against $g(y)\dd y$ gives the values $0$
and $\Circ_g(r)/(2\pi r)$; Fubini's theorem applies because
$\oint_{C_r}|\KBS(x-y)|\dd s(x)$ is bounded by a constant, depending on $r$, times
$1+|\log\operatorname{dist}(y,C_r)|$, which is integrable in $y$ over compact sets,
and $\{|y|=r\}$ has measure $0$. Hence $\KBS\ast g=\Circ_g(|x|)\,x^\perp/(2\pi|x|^2)$.
With $g=\Phi$ and the definition of $h$, on $[0,T]\times U$,
\[
v=\KBS\ast\Phi+\KBS\ast(\omega_{\mathrm r}-\Phi)+\KBS\ast(\omega_{\mathrm i}+\omega_{\mathrm o})
=\Big(\frac{\Circ_\Phi(|x|)}{2\pi|x|^2}+\frac G{|x|^2}\Big)x^\perp+\KBS\ast(\omega_{\mathrm r}-\Phi)+h,
\]
and the bracket is $\Omega(|x|;G)$. Since $\partial_\tau\flow_\tau(y)=v(\tau,\flow_\tau(y))$
and $x^\perp\cdot\er=0$, we get $\frac{\dd}{\dd\tau}|\flow_\tau(y)|=
\big((h+\KBS\ast(\omega_{\mathrm r}-\Phi))\cdot\er\big)(\tau,\flow_\tau(y))$ while
$\flow_\tau(y)\in N'$, which gives the integral formula.

\emph{Proof of (3).} Since $\supp\Phi\subset\{a\le|y|\le1\}\subset N$
(Lemma~I.\ref{I-lem:ring-facts}), $w(t)$ is a smooth, real, $3$-fold rotationally symmetric function
supported in $N$. The derivatives $\partial^\alpha_xw$ and
$\partial_t\partial^\alpha_xw$ are uniformly continuous on the compact set
$[0,T]\times N$ and vanish outside it, so the difference quotients of
$t\mapsto w(t)$ converge in every $H^k$, and $w\in C^1([0,T];X^k)$.
Differentiating $\omega_{\mathrm r}(t,\flow_t(y))=\omega_{\mathrm r}(0,y)$ in $t$ gives
$\partial_t\omega_{\mathrm r}+v\cdot\nabla\omega_{\mathrm r}=0$. On $N$ we insert (2),
$x^\perp\cdot\nabla\Phi=\partial_\theta\Phi=0$ and $x^\perp\cdot\nabla w=\partial_\theta w$:
\[
\partial_tw=-\Omega(\cdot;G)\,\partial_\theta w-(\KBS\ast w)\cdot\nabla\Phi-(\KBS\ast w)\cdot\nabla w-h\cdot\nabla w-h\cdot\nabla\Phi .
\]
The first two terms are $\Lin(G)w$ by Definition~I.\ref{I-def:sector-setting}, and the
third is $B(w,w)$. Outside $N$ both sides of Equation~\eqref{eq:ring} vanish.
\end{proof}

In the setting of Lemma~\ref{lem:reduction}, put $w=\omega_{\mathrm r}-\Phi$. Then
$t\mapsto\|w(t)\|_{H^k}$ is continuous
on $[0,T]$ for every $k$, since $\omega_{\mathrm r}$ is smooth and its support stays
in a fixed compact set (Theorem~\ref{thm:smooth}\textup{(1)}). We fix once and for
all the auxiliary objects of the estimates below. The functions we estimate are
supported in $N$, and the velocities that act on them are needed only on a
neighborhood of $N$. We therefore multiply the Biot-Savart velocity by a compactly
supported cutoff $\chi_B$ that does not change it on $N'$, so that
Lemma~\ref{lem:Hk-calculus}\textup{(3)} applies, and we use a cutoff $\chi_N$ around $N$ in Lemma~\ref{lem:Hk-calculus}\textup{(1)}
and \textup{(2)}. Let $N_1=\{\frac a4<|y|<3\}$ be
the interior of $N'$. Fix $\chi_B\in C_c^\infty(B_4)$ with $\chi_B=1$ on $B_{7/2}$,
and $\chi_N\in C_c^\infty(N_1)$ with $\chi_N=1$ on a neighborhood of $N$. For each
integer $j\ge0$ fix a constant $C^{\mathrm{el}}_j$ of
Lemma~\ref{lem:Hk-calculus}\textup{(3)} for $j$, $\chi_B$ and $\rho=4$; it depends
only on $j$. For each integer $k\ge2$ fix a constant $C_P(k)$ of
Lemma~\ref{lem:Hk-calculus} for $\max(k,3)$, the compact set $N$, the open set
$N_1$ and $\chi_N$; it depends only on $k$ and $a$.

\begin{lemma}
\label{lem:bilinear}
Let $k\ge2$ be an integer. Put $C_B=C_B(k)=2C_P(k)C^{\mathrm{el}}_{k-1}$ and
$C_{\mathrm{BS}}=2C^{\mathrm{el}}_1$, and let $C_h=C_h(k)$ be the supremum of
$\|h\cdot\nabla w\|_{H^k}$ over all $w\in X^{k+1}$ and
$h\in C^k(N';\R^2)$ with $\|w\|_{H^{k+1}}\le1$ and $\|h\|_{C^k(N')}\le1$. Then
$C_h$ is finite and depends only on $k$ and $a$, and for $v,w\in X^{k+1}$ and
every vector field $h\in C^k(N';\R^2)$:
\begin{enumerate}[label=\textup{(\arabic*)}]
\item $\|B(v,w)\|_{H^k}\le C_B\|v\|_{H^k}\|w\|_{H^{k+1}}$;
\item $\|h\cdot\nabla w\|_{H^k}\le C_h\|h\|_{C^k(N')}\|w\|_{H^{k+1}}$;
\item $\|h\cdot\nabla\Phi\|_{H^k}\le\Cnab\|h\|_{C^k(N')}$, with $\Cnab=\Cnab(k)$ of
  Lemma~I.\ref{I-lem:uniform};
\item $\sup_{|x|\le3}|(\KBS\ast v)(x)|\le C_{\mathrm{BS}}\|v\|_{H^1}$.
\end{enumerate}
The functions in \textup{(1)}--\textup{(3)} are supported in $N$, and they lie in
$X^k$ if $h$ is $3$-fold rotationally symmetric.
\end{lemma}
\begin{proof}
Since $\nabla w$ is supported in $N$,
$B(v,w)=-\sum_{i=1,2}(\chi_B\KBS\ast v)_i\,\partial_iw$. By
Lemma~\ref{lem:Hk-calculus}\textup{(2)} (with $\max(k,3)$ in place of its $k$, and
$j=k$) and \textup{(3)},
\[
\|B(v,w)\|_{H^k}\le2C_P(k)\|\chi_B\KBS\ast v\|_{H^k}\max_i\|\partial_iw\|_{H^k}\le2C_P(k)C^{\mathrm{el}}_{k-1}\|v\|_{H^{k-1}}\|w\|_{H^{k+1}},
\]
which is (1). Next, $h\cdot\nabla w=(\chi_Nh)\cdot\nabla w$, and the Leibniz rule
$\partial^\alpha\big((\chi_Nh)\cdot\nabla w\big)=\sum_{\beta\le\alpha}\binom\alpha\beta
\partial^\beta(\chi_Nh)\cdot\partial^{\alpha-\beta}\nabla w$ with
$|\partial^\beta(\chi_Nh)|\le2^{|\beta|}\|\chi_N\|_{C^k}\|h\|_{C^k(N')}$ bounds
$\|h\cdot\nabla w\|_{H^k}$ by a constant depending only on $k$ and $\chi_N$, hence
on $k$ and $a$, times $\|h\|_{C^k(N')}\|w\|_{H^{k+1}}$. So $C_h$ is finite, depends
only on $k$ and $a$, and (2) holds by homogeneity in $h$ and in $w$. Item (3) is the
last inequality of Lemma~I.\ref{I-lem:uniform}. For (4), $\KBS\ast v=\chi_B\KBS\ast v$ on
$B_{7/2}$, and Lemma~\ref{lem:embed}\textup{(E2)}, applied to each component, and
Lemma~\ref{lem:Hk-calculus}\textup{(3)} give
$\|\chi_B\KBS\ast v\|_{L^\infty}\le2C^{\mathrm{el}}_1\|v\|_{H^1}=C_{\mathrm{BS}}\|v\|_{H^1}$,
which gives (4), since $\{|x|\le3\}\subset B_{7/2}$. The functions in
(1)--(3) vanish outside $\supp\nabla w\cup\supp\nabla\Phi\subset N$. If $v$, $w$ and
$h$ are $3$-fold rotationally symmetric, so are the vector fields $\KBS\ast v$ (by the
equivariance in the proof of Lemma~\ref{lem:reduction}), $\nabla w$ and
$\nabla\Phi$, hence the three scalar products.
\end{proof}

Lemma~\ref{lem:energy} below compares the deviation $w$ of a ring with an
approximate solution $w_{\mathrm{app}}$ by an energy estimate in $H^s$. Its rate
$\Lam(s)$, defined next, bounds the transport of $w-w_{\mathrm{app}}$ by the
velocity and the coupling through $\nabla(\Phi+w_{\mathrm{app}})$; the entry
$2\max_{G\in[0,\Gone]}\sigma(G)$ in its definition makes $\Lam(s)\ge2\sigma(G)$ for every $G$,
which the proof of Lemma~\ref{lem:window} uses.
This rate involves only norms on $N'$. Let $s\ge3$ be
an integer, let $n_s$ be the number of multi-indices $\alpha\in\N_0^2$ with
$|\alpha|\le s$, and put $C_{\mathrm{HK}}=C_P(s)$. The function
$\Omega(r;G)=\Omega(r;0)+Gr^{-2}$ is smooth on $[\frac a4,3]$ and affine in $G$, so
\[
C_\Omega:=\sup_{G\in[0,\Gone]}\Big(\|\Omega(|x|;G)x^\perp\|_{W^{1,\infty}(N_1)}+\|\Omega(|x|;G)x^\perp\|_{H^s(N_1)}\Big)<\infty .
\]
Let $C_{N'}$ be the least constant such that
$\|h\|_{W^{1,\infty}(N_1)}+\|h\|_{H^s(N_1)}\le C_{N'}\|h\|_{C^s(N')}$ for every
$h\in C^s(N';\R^2)$; it is finite and depends only on $s$ and $a$. Put
$A_s=C_\Omega+C_{N'}+C^{\mathrm{el}}_s+6C^{\mathrm{el}}_2$ and
\begin{equation}
\label{eq:Lambda}
\Lam(s)=\max\Big(2\max_{G\in[0,\Gone]}\sigma(G),\ n_sC_{\mathrm{HK}}A_s+2C_{\mathrm{HK}}C^{\mathrm{el}}_{s-1}\big(\|\Phi\|_{H^{s+1}}+1\big)\Big).
\end{equation}

\begin{lemma}
\label{lem:energy}
Let $s\ge3$ be an integer, and let $\Lam(s)$ be the number of
Equation~\eqref{eq:Lambda}; it satisfies $\Lam(s)\ge2\max_{G\in[0,\Gone]}\sigma(G)$,
where the maximum exists by Proposition~I.\ref{I-prop:U}\textup{(a)}, and it depends
only on $\Phi$, $\Gone$ and $s$. Let $G\in[0,\Gone]$
and $t>0$, and let the assumptions of Lemma~\ref{lem:reduction}\textup{(3)} hold on
$[0,t]$, with $w=\omega_{\mathrm r}-\Phi$. Let
$w_{\mathrm{app}}\in C([0,t];X^{s+1})\cap C^1([0,t];X^s)$ be real-valued, and put
$R_{\mathrm{app}}=\partial_tw_{\mathrm{app}}-\big(\Lin(G)w_{\mathrm{app}}+B(w_{\mathrm{app}},w_{\mathrm{app}})-h\cdot\nabla w_{\mathrm{app}}-h\cdot\nabla\Phi\big)$
and $v=w-w_{\mathrm{app}}$. If on $[0,t]$ we have $\|w\|_{H^s}\le1$,
$\|w_{\mathrm{app}}\|_{H^{s+1}}\le1$ and $\|h\|_{C^{s+1}(N')}\le1$, then for
$\tau\in[0,t]$,
\[
\|v(\tau)\|_{H^s}\le e^{\Lam(s)\tau}\|v(0)\|_{H^s}+\int_0^\tau e^{\Lam(s)(\tau-\tau')}\|R_{\mathrm{app}}(\tau')\|_{H^s}\dd\tau' .
\]
\end{lemma}
\begin{proof}
On $N_1$ put $U=\Omega(|x|;G)x^\perp+h+\KBS\ast w$, and for $f$ supported in $N$ let
$U\cdot\nabla f$ mean $(\chi_NU)\cdot\nabla f$. By
Lemma~\ref{lem:reduction}\textup{(3)},
$\partial_tw+U\cdot\nabla w=-(\KBS\ast w)\cdot\nabla\Phi-h\cdot\nabla\Phi$, and by
the definition of $R_{\mathrm{app}}$ the same identity holds for $w_{\mathrm{app}}$,
with $\KBS\ast w_{\mathrm{app}}$ in place of $\KBS\ast w$ in $U$ and with
$R_{\mathrm{app}}$ added on the right. Since
$(\KBS\ast w)\cdot\nabla w-(\KBS\ast w_{\mathrm{app}})\cdot\nabla w_{\mathrm{app}}=(\KBS\ast w)\cdot\nabla v+(\KBS\ast v)\cdot\nabla w_{\mathrm{app}}$,
subtraction gives
\[
\partial_tv+U\cdot\nabla v=-(\KBS\ast v)\cdot\nabla(\Phi+w_{\mathrm{app}})-R_{\mathrm{app}},
\]
where every term lies in $C([0,t];H^s)$, and
$v\in C^1([0,t];H^s)\cap C([0,t];H^{s+1})$ is real and supported in $N$.

The field $U$ is divergence-free on $N_1$: $\divop h=0$ by
Lemma~\ref{lem:reduction}\textup{(1)}, $\divop(\KBS\ast w)=0$, and
$\divop(f(|x|)x^\perp)=f'(|x|)\frac{x}{|x|}\cdot x^\perp=0$. By
Lemma~\ref{lem:Hk-calculus}\textup{(3)} and Lemma~\ref{lem:embed}\textup{(E2)},
applied to each entry, and $\|w\|_{H^s}\le1$,
$\|\KBS\ast w\|_{H^s(N_1)}\le\|\chi_B\KBS\ast w\|_{H^{s+1}}\le C^{\mathrm{el}}_s$ and
$\|\KBS\ast w\|_{W^{1,\infty}(N_1)}\le6\|\chi_B\KBS\ast w\|_{H^3}\le6C^{\mathrm{el}}_2$.
With the definitions of $C_\Omega$ and $C_{N'}$ and $\|h\|_{C^s(N')}\le1$, this gives
$\|U\|_{W^{1,\infty}(N_1)}+\|U\|_{H^s(N_1)}\le A_s$. Since $U$ and $v$ are real and
$\divop U=0$ on $N_1$, Lemma~\ref{lem:Hk-calculus}\textup{(1)} gives
$|\langle\partial^\alpha(U\cdot\nabla v),\partial^\alpha v\rangle|\le C_{\mathrm{HK}}A_s\|v\|_{H^s}^2$
for $|\alpha|\le s$. Since $\nabla(\Phi+w_{\mathrm{app}})$ is supported in $N$,
Lemma~\ref{lem:Hk-calculus}\textup{(2)} and \textup{(3)} give
\[
\begin{aligned}
\|(\KBS\ast v)\cdot\nabla(\Phi+w_{\mathrm{app}})\|_{H^s}
&\le2C_{\mathrm{HK}}\|\chi_B\KBS\ast v\|_{H^s}\big(\|\Phi\|_{H^{s+1}}+\|w_{\mathrm{app}}\|_{H^{s+1}}\big)\\
&\le2C_{\mathrm{HK}}C^{\mathrm{el}}_{s-1}\big(\|\Phi\|_{H^{s+1}}+1\big)\|v\|_{H^s}.
\end{aligned}
\]
We apply $\partial^\alpha$ to the equation of $v$, take the $L^2$ inner product with
$\partial^\alpha v$, and sum over $|\alpha|\le s$. The Cauchy-Schwarz inequality and
Equation~\eqref{eq:Lambda} give
$\frac12\frac{\dd}{\dd\tau}\|v\|_{H^s}^2\le\Lam(s)\|v\|_{H^s}^2+\|R_{\mathrm{app}}\|_{H^s}\|v\|_{H^s}$.
To divide by $\|v\|_{H^s}$, which may vanish, we regularize it. For $\epsilon>0$
the function $y_\epsilon=(\|v\|_{H^s}^2+\epsilon^2)^{1/2}$ is $C^1$, as
$v\in C^1([0,t];H^s)$ and $y_\epsilon\ge\epsilon$. Since
$y_\epsilon y_\epsilon'=\frac12\frac{\dd}{\dd\tau}\|v\|_{H^s}^2$ and
$\|v\|_{H^s}\le y_\epsilon$, it satisfies
$y_\epsilon'\le\Lam(s)y_\epsilon+\|R_{\mathrm{app}}\|_{H^s}$, so
$(e^{-\Lam(s)\tau}y_\epsilon)'\le e^{-\Lam(s)\tau}\|R_{\mathrm{app}}\|_{H^s}$.
Integrating from $0$ to $\tau$, and letting $\epsilon\to0$ with
$\|v\|_{H^s}\le y_\epsilon$ and $y_\epsilon(0)\to\|v(0)\|_{H^s}$, proves the lemma.
\end{proof}

\emph{Approximate solutions.} Following Grenier, we approximate the deviation $w$
by the first $J$ terms $w^{[1]},\dots,w^{[J]}$ of its formal expansion in powers of
the size of the data, in which $h$ and the initial value count as terms of first
order: $w^{[1]}$ solves the linear equation forced by $-h\cdot\nabla\Phi$, with an
initial value $w_0$ that is $0$ or $w(0)$ in our applications, and $w^{[j]}$,
$j\ge2$, solves the linear equation with initial value $0$, forced by the terms of
order $j$ of $B(w,w)-h\cdot\nabla w$. The
residual of the sum then consists of terms of order larger than $J$
(Equation~\eqref{eq:hier-res}), and Lemma~\ref{lem:energy} controls the difference
between $w$ and the sum. Let the assumptions of Lemma~\ref{lem:reduction} hold
on $[0,T]$, assume $G\in[0,\Gone]$, and put $\Lin=\Lin(G)$ and $F=-h\cdot\nabla\Phi$.
By Lemma~\ref{lem:reduction}\textup{(1)} and Lemma~\ref{lem:bilinear}, $F$ is real
and lies in $C([0,T];X^k)$ for every $k$; so, by Lemma~\ref{lem:duhamel}, for every
real $w_0\in\bigcap_kX^k$ the function
$e^{t\Lin}w_0+\int_0^te^{(t-\tau)\Lin}F(\tau)\dd\tau$ is real, lies in
$C([0,T];X^k)\cap C^1([0,T];X^{k-1})$ for every $k\ge1$, and solves
$\partial_tu=\Lin u+F$. If $\supp\omega_{\mathrm r}(0)\subset N$, this applies to
$w_0=w(0)=\omega_{\mathrm r}(0)-\Phi$, which is then smooth
(Lemma~\ref{lem:reduction}\textup{(2)}), real, $3$-fold rotationally symmetric and supported in $N$,
hence lies in every $X^k$. Let $w^{[1]}$ be such a function, let $J\ge1$ be an integer,
and for $2\le j\le J$ define recursively
\begin{equation}
\label{eq:hier}
w^{[j]}(t)=\int_0^te^{(t-\tau)\Lin}g^{[j]}(\tau)\dd\tau,\qquad
g^{[j]}=\sum_{i+l=j}B(w^{[i]},w^{[l]})-h\cdot\nabla w^{[j-1]},
\end{equation}
where here and below $i,l\ge1$ in the sums. By induction on $j$,
Lemma~\ref{lem:bilinear} (whose bounds, with the continuity of $t\mapsto h(t)$ in
every $C^k(N')$, give continuity in time) and Lemma~\ref{lem:duhamel} show that
$g^{[j]}$ and $w^{[j]}$ are real, $g^{[j]}\in C([0,T];X^k)$ and
$w^{[j]}\in C([0,T];X^k)\cap C^1([0,T];X^{k-1})$ for every $k\ge1$, and
$\partial_tw^{[j]}=\Lin w^{[j]}+g^{[j]}$. Put $w_{\mathrm{app}}=\sum_{j=1}^Jw^{[j]}$.
Summing the equations of the $w^{[j]}$, and using $F+h\cdot\nabla\Phi=0$,
$\sum_{j=2}^{J}\sum_{i+l=j}B(w^{[i]},w^{[l]})=B(w_{\mathrm{app}},w_{\mathrm{app}})-\sum_{i,l\le J,\,i+l>J}B(w^{[i]},w^{[l]})$
and $\sum_{j=2}^{J}h\cdot\nabla w^{[j-1]}=h\cdot\nabla w_{\mathrm{app}}-h\cdot\nabla w^{[J]}$,
we find that the residual of $w_{\mathrm{app}}$ in Lemma~\ref{lem:energy} is
\begin{equation}
\label{eq:hier-res}
R_{\mathrm{app}}=-\sum_{i,l\le J,\,i+l>J}B(w^{[i]},w^{[l]})+h\cdot\nabla w^{[J]} .
\end{equation}

\subsection{A Quiet Ring}

A ring is quiet while the instability has not reached it: it starts exactly at
$\Phi$, so its deviation is created by the field $h$ alone, and
Theorem~\ref{thm:quiet} shows that it stays small. We bound a norm of the field by a
function $\eta$ of time, and its effect by the Duhamel majorant of $\eta$, defined
next. Since the semigroup bound of Proposition~I.\ref{I-prop:semigroup} grows like
$e^{\sigma(G)t}$, the majorant bounds, up to a constant factor, the linear
response to a forcing of size $\eta$.

\begin{definition}
\label{def:majorant}
Let $G\in[0,\Gone]$, $T>0$, and let $\eta\colon[0,T]\to[0,\infty)$ be integrable.
The \textbf{Duhamel majorant} of $\eta$ is
$\maj(t)=\int_0^te^{\sigma(G)(t-\tau)}\eta(\tau)\dd\tau$, $t\in[0,T]$.
\end{definition}

\begin{lemma}
\label{lem:majorant}
Let $G$, $T$, $\eta$ and $\maj$ be as in Definition~\ref{def:majorant}. Then for
$0\le\tau\le t\le T$:
\textup{(M1)} $e^{\sigma(G)(t-\tau)}\maj(\tau)\le\maj(t)$;
\textup{(M2)} $\int_\tau^t\eta\le\maj(t)$;
\textup{(M3)} $\int_0^t\maj\le\maj(t)/\sigma(G)$.
\end{lemma}
\begin{proof}
Write $\sigma=\sigma(G)$, which is positive by Proposition~I.\ref{I-prop:SU}. We have
$\maj(t)=e^{\sigma(t-\tau)}\maj(\tau)+\int_\tau^te^{\sigma(t-v)}\eta(v)\dd v$, and the
integral is nonnegative, which gives (M1), and at least $\int_\tau^t\eta$, which
gives (M2). By (M1), $\maj(\tau)\le e^{-\sigma(t-\tau)}\maj(t)$; integrating over
$\tau\in[0,t]$ gives (M3).
\end{proof}

\begin{theorem}
\label{thm:quiet}
Assume \textup{(N1)}--\textup{(N3)}. Let $s_h\ge3$ and $N_Q\ge2$ be integers with
$N_Q\sigb\ge\Lam(s_h)$, and put $M_h=s_h+N_Q+1$. There are $\majst\in(0,1]$ and
$C_Q<\infty$, depending only on $(\Phi,\Gone,s_h,N_Q)$, such that the following
holds. Let $\omega$,
$\omega_{\mathrm r},\omega_{\mathrm i},\omega_{\mathrm o}$, $G$ and $h$ be as in
Lemma~\ref{lem:reduction} on $[0,T]$, with $\omega_{\mathrm r}(0)=\Phi$ and
$G\in[0,\Gone]$, and put $w=\omega_{\mathrm r}-\Phi$. Let $\eta$ be measurable with
$\|h(t)\|_{C^{M_h}(N')}\le\eta(t)\le1$ on $[0,T]$, let $\maj$ be its Duhamel
majorant, and assume $\maj\le\majst$ on $[0,T]$. Then for every $t\in[0,T]$:
\begin{enumerate}[label=\textup{(\alph*)}]
\item $\supp\omega_{\mathrm r}(t)\subset\NNm$;
\item $\|w(t)-\wlin(t)\|_{H^{s_h}}\le C_Q\maj(t)^2$, where
  $\wlin(t)=\int_0^te^{(t-\tau)\Lin(G)}\big(-h(\tau)\cdot\nabla\Phi\big)\dd\tau$
  is the \textbf{linear response};
\item $\|w(t)\|_{H^{s_h}}\le C_Q\maj(t)$.
\end{enumerate}
\end{theorem}
\begin{proof}
Throughout, $\sigma=\sigma(G)\ge\sigb$ and $\Lin=\Lin(G)$. We compare $w$ with the
approximate solution of Equation~\eqref{eq:hier} with $J=N_Q$ and $w^{[1]}=\wlin$,
bound its terms by powers of $\maj$, and close with Lemma~\ref{lem:energy} in a
continuity argument on the radial displacement and on $\|w\|_{H^{s_h}}$. Both
terms of the bound on the residual contain the factor $\maj^{N_Q}$
(Equation~\eqref{eq:Q-res} below), and the condition $N_Q\sigb\ge\Lam(s_h)$ lets
the growth of this factor absorb the growth $e^{\Lam(s_h)t}$ allowed by the
energy estimate.

\emph{Constants.} Put $\sigma'_0=\frac12(\frac9{200}+\sigb)\in(\frac9{200},\sigb)$ and
\[
C_L=\max_{0\le k\le M_h}C_S(k,\sigma'_0),\qquad
C_{\mathrm{bil}}=\max_{2\le k\le M_h}\max\big(C_B(k),C_h(k),\Cnab(k)\big),
\]
with $C_S$ of Proposition~I.\ref{I-prop:semigroup} and the constants of
Lemma~\ref{lem:bilinear}; so $\|e^{t\Lin}\|_{X^k\to X^k}\le C_Le^{\sigma t}$ for
$k\le M_h$ and $t\ge0$. For $1\le j\le N_Q$ put $k_j=M_h-j$, so that
$k_{N_Q}=s_h+1$ and $k_j+1=k_{j-1}$. Put $C^{(1)}=C_LC_{\mathrm{bil}}$ and, for
$2\le j\le N_Q$,
\[
\begin{gathered}
C^{(j)}=C_LC_{\mathrm{bil}}\Big(\frac1{(j-1)\sigb}\sum_{i+l=j}C^{(i)}C^{(l)}+C^{(j-1)}\Big),\\
C_{\mathrm{res}}=C_{\mathrm{bil}}\Big(\sum_{i,l\le N_Q,\,i+l>N_Q}C^{(i)}C^{(l)}+C^{(N_Q)}\Big),\\
C_Q=\sum_{j=1}^{N_Q}C^{(j)}+C_{\mathrm{res}}\big(1+\sigb^{-1}\big),\qquad
C_\Delta=1+\frac{C_{\mathrm{BS}}C_Q}{\sigb},\\
\majst=\min\Big(1,\frac1{2C_Q},\frac a{16C_\Delta}\Big).
\end{gathered}
\]
The names $C_{\mathrm{bil}}$, $C^{(j)}$ and $C_{\mathrm{res}}$ are local to this proof;
later in this section they denote other numbers.
These numbers depend only on $(\Phi,\Gone,s_h,N_Q)$, and $C_\Delta\majst\le a/16$.

\emph{Approximate solution.} Let $w^{[j]}$, $1\le j\le N_Q$, be the functions of
Equation~\eqref{eq:hier} with $J=N_Q$ and $w^{[1]}=\wlin$ (the case $w_0=0$), so that
they are regular by Lemmas~\ref{lem:duhamel} and \ref{lem:bilinear}. By
Lemma~\ref{lem:bilinear}\textup{(3)},
$\|F(\tau)\|_{H^{M_h}}\le C_{\mathrm{bil}}\eta(\tau)$. We claim that
\begin{equation}
\label{eq:Q-hier}
\|w^{[j]}(t)\|_{H^{k_j}}\le C^{(j)}\maj(t)^j\qquad(t\in[0,T],\ 1\le j\le N_Q).
\end{equation}
For $j=1$,
$\|\wlin(t)\|_{H^{k_1}}\le C_L\int_0^te^{\sigma(t-\tau)}C_{\mathrm{bil}}\eta(\tau)\dd\tau=C^{(1)}\maj(t)$.
Let $j\ge2$ and assume Equation~\eqref{eq:Q-hier} for all indices less than $j$.
Since $k_j+1=k_{j-1}\le k_l$ for $l\le j-1$, Lemma~\ref{lem:bilinear} gives
$\|g^{[j]}(\tau)\|_{H^{k_j}}\le C_{\mathrm{bil}}\big(\sum_{i+l=j}C^{(i)}C^{(l)}\maj(\tau)^j+C^{(j-1)}\eta(\tau)\maj(\tau)^{j-1}\big)$.
By (M1), $e^{\sigma(t-\tau)}\maj(\tau)^j\le e^{-(j-1)\sigma(t-\tau)}\maj(t)^j$, whose
integral over $[0,t]$ is at most $\maj(t)^j/((j-1)\sigb)$; and
$e^{\sigma(t-\tau)}\eta(\tau)\maj(\tau)^{j-1}\le\eta(\tau)\maj(t)^{j-1}$, whose integral
is at most $\maj(t)^j$ by (M2). With $\|e^{(t-\tau)\Lin}\|\le C_Le^{\sigma(t-\tau)}$ this
gives Equation~\eqref{eq:Q-hier} for $j$. Since $k_j\ge s_h+1$ and $\maj\le1$,
Equations~\eqref{eq:hier-res} and \eqref{eq:Q-hier} and Lemma~\ref{lem:bilinear} give
\begin{equation}
\label{eq:Q-res}
\|R_{\mathrm{app}}(\tau)\|_{H^{s_h}}\le C_{\mathrm{res}}\big(\maj(\tau)^{N_Q+1}+\eta(\tau)\maj(\tau)^{N_Q}\big),
\end{equation}
and $\|w_{\mathrm{app}}\|_{H^{s_h+1}}\le\sum_jC^{(j)}\maj\le C_Q\majst\le\frac12$ on
$[0,T]$.

\emph{Continuity argument.} Let $t_*$ be the largest $t\in[0,T]$ such that
$\delta_{\mathrm r}(t)\le a/8$ and $\|w(\tau)\|_{H^{s_h}}\le1$ for $\tau\in[0,t]$; it
exists and is positive since $\delta_{\mathrm r}(0)=0$, $w(0)=0$, $\delta_{\mathrm r}$
is continuous (see its definition at the start of this section), and
$\|w\|_{H^{s_h}}$ is continuous by Theorem~\ref{thm:smooth}\textup{(1)}. Because $\supp\Phi\subset\{a\le|y|\le1\}$ and
$1+\frac a8\le\frac54$, the support of $\omega_{\mathrm r}(\tau)$ lies in
$\NNm\subset N$ for $\tau\le t_*$. So Lemma~\ref{lem:reduction}\textup{(3)} and
Lemma~\ref{lem:energy} with $s=s_h$ apply on $[0,t_*]$; the conditions
$\|w_{\mathrm{app}}\|_{H^{s_h+1}}\le1$ and $\|h\|_{C^{s_h+1}(N')}\le\eta\le1$ hold.
Since $v=w-w_{\mathrm{app}}$ has $v(0)=0$, Equation~\eqref{eq:Q-res} gives, for
$t\le t_*$,
\[
\|v(t)\|_{H^{s_h}}\le C_{\mathrm{res}}\int_0^te^{\Lam(s_h)(t-\tau)}\maj(\tau)^{N_Q}\big(\maj(\tau)+\eta(\tau)\big)\dd\tau .
\]
For $0\le\tau\le t$, (M1) and $N_Q\sigma\ge N_Q\sigb\ge\Lam(s_h)$ give
$e^{\Lam(s_h)(t-\tau)}\maj(\tau)^{N_Q}\le e^{(\Lam(s_h)-N_Q\sigma)(t-\tau)}\maj(t)^{N_Q}$,
which is at most $\maj(t)^{N_Q}$. By (M2) and (M3) we conclude, for $t\le t_*$,
\begin{equation}
\label{eq:Q-v}
\|v(t)\|_{H^{s_h}}\le C_{\mathrm{res}}\maj(t)^{N_Q}\int_0^t(\maj+\eta)\le C_{\mathrm{res}}\big(1+\sigb^{-1}\big)\maj(t)^{N_Q+1}.
\end{equation}
Since $w=\wlin+\sum_{j\ge2}w^{[j]}+v$ and $\maj\le1$, Equations~\eqref{eq:Q-hier}
and \eqref{eq:Q-v} give, for $t\le t_*$,
\begin{equation}
\label{eq:Q-b}
\|w(t)-\wlin(t)\|_{H^{s_h}}\le C_Q\maj(t)^2,\qquad\|w(t)\|_{H^{s_h}}\le C_Q\maj(t)\le\tfrac12 .
\end{equation}
For $y\in\supp\Phi$ and $\tau\le t\le t_*$, the point $\flow_\tau(y)$ lies in $N$, so
by Lemma~\ref{lem:reduction}\textup{(2)}, Lemma~\ref{lem:bilinear}\textup{(4)},
Equation~\eqref{eq:Q-b}, (M2) and (M3),
\[
\big||\flow_\tau(y)|-|y|\big|\le\int_0^t\big(\eta+C_{\mathrm{BS}}\|w\|_{H^{s_h}}\big)\le\maj(t)+\frac{C_{\mathrm{BS}}C_Q}{\sigb}\maj(t)=C_\Delta\maj(t).
\]
So $\delta_{\mathrm r}\le C_\Delta\maj\le a/16$ and $\|w\|_{H^{s_h}}\le\frac12$ on
$[0,t_*]$. If $t_*<T$, continuity would give a larger $t$ with
$\delta_{\mathrm r}(t)\le a/8$ and $\|w\|_{H^{s_h}}\le1$ on $[0,t]$; hence $t_*=T$.
Then (a) follows from $\delta_{\mathrm r}\le C_\Delta\maj$ and Lemma~\ref{lem:annuli},
and (b), (c) are
Equation~\eqref{eq:Q-b}.
\end{proof}

\subsection{An Active Ring}

A ring is active once its deviation from $\Phi$ has a component along the
unstable eigenfunction that grows at the rate $\sigma(G)$. Theorem~\ref{thm:active}
shows that the deviation then follows this component, up to a decaying transient
and errors of relative size $\hat a$, $\kappa^{1/2}$ and $\eta_b$, until its
envelope $\kappa e^{\sigma(G)u}$ reaches a small size $\delta_0$.
Corollary~\ref{cor:active-moment} gives the outer moment of the ring, through
which it drives the rings outside it. The unstable component of the linear
response is explicit, by Proposition~I.\ref{I-prop:semigroup}; the next lemma
writes it with a running coefficient $c(u)$, which accounts for the forcing by $h$.

\begin{lemma}
\label{lem:running}
Let $G\in[0,\Gone]$, $\bar u>0$, $w_0\in X^0$ real, and let $h$ be continuous on
$[0,\bar u]$ with values in $C^0(N';\R^2)$ and $3$-fold rotationally symmetric. Put
$F(u)=-h(u)\cdot\nabla\Phi$ and
\begin{equation}
\label{eq:running}
c(u)=\ellst(w_0)+\int_0^ue^{-\lamst v}\,\ellst(F(v))\dd v\qquad(0\le u\le\bar u),
\end{equation}
the \textbf{running coefficient}. Then
$\wlin_0(u):=e^{u\Lin(G)}w_0+\int_0^ue^{(u-v)\Lin(G)}F(v)\dd v$ satisfies
$\ellst(\wlin_0(u))=c(u)e^{\lamst u}$ and
$\Riesz\,\wlin_0(u)=2\Re\big(c(u)e^{\lamst u}\phist\big)$ for $0\le u\le\bar u$.
\end{lemma}
\begin{proof}
The function $F$ is continuous on $[0,\bar u]$ with values in $X^0$, since
$\|h(u)\cdot\nabla\Phi-h(u')\cdot\nabla\Phi\|_{L^2}\le\|h(u)-h(u')\|_{C^0(N')}\|\nabla\Phi\|_{L^2}$,
and $F(u)$ is real, supported in $N$ and $3$-fold rotationally symmetric. The functional
$\ellst$ is bounded on $X^0$, so it commutes with the integral, and
Proposition~I.\ref{I-prop:semigroup} gives $\ellst(e^{\tau\Lin}f)=e^{\lamst\tau}\ellst(f)$
for $f\in X^0$. Hence
$\ellst(\wlin_0(u))=e^{\lamst u}\ellst(w_0)+\int_0^ue^{\lamst(u-v)}\ellst(F(v))\dd v=c(u)e^{\lamst u}$.
Since $\wlin_0(u)$ is real, Definition~I.\ref{I-def:eigendata} gives
$\Riesz\,\wlin_0(u)=2\Re\big(\ellst(\wlin_0(u))\phist\big)=2\Re\big(c(u)e^{\lamst u}\phist\big)$.
\end{proof}

On the range of $1-\Riesz$ the semigroup grows at most at a rate $\sigma'<\sigb$
(Proposition~I.\ref{I-prop:semigroup}). The next lemma uses this to bound the rest
$(1-\Riesz)\wlin_0$ of the linear response.

\begin{lemma}
\label{lem:linear-part}
Assume \textup{(N1)}--\textup{(N3)}. Let $\sigma'\in(\frac9{200},\sigb)$, and let
$p\ge0$ and $k\ge2$ be integers. Put
$C_{\mathrm{lin}}(k)=C_S(k,\sigma')\big(1+\Cnab(k)/\sigma'+\Cnab(k)/(\sigb-\sigma')\big)$,
with $C_S$ of Proposition~I.\ref{I-prop:semigroup}. Let $G\in[0,\Gone]$,
$\sigma=\sigma(G)$, $\lambda=\lamst(G)$, $\bar u>0$, $A\ge0$, $\kappa\ge0$ and
$\eta_b\ge0$. Let $w_0\in X^k$ be real with $\|w_0\|_{H^k}\le A$, and let $h$ be
continuous on $[0,\bar u]$ with values in $C^k(N';\R^2)$, $3$-fold rotationally symmetric, with
$\|h(u)\|_{C^k(N')}\le A\jap{u}^p+\eta_b\kappa e^{\sigma u}$. Let $c$ and $\wlin_0$ be as
in Lemma~\ref{lem:running}. Then for $u\in[0,\bar u]$,
\[
\begin{gathered}
\wlin_0(u)=2\Re\big(c(u)e^{\lambda u}\phist\big)+(1-\Riesz)\wlin_0(u),\\
\|(1-\Riesz)\wlin_0(u)\|_{H^k}\le C_{\mathrm{lin}}(k)\big(A\jap{u}^pe^{\sigma'u}+\eta_b\kappa e^{\sigma u}\big).
\end{gathered}
\]
\end{lemma}
\begin{proof}
The decomposition is Lemma~\ref{lem:running}. For real $f\in X^k$,
Proposition~I.\ref{I-prop:semigroup} gives
$\Riesz e^{\tau\Lin}f=2\Re\big(\ellst(e^{\tau\Lin}f)\phist\big)=2\Re\big(e^{\lambda\tau}\ellst(f)\phist\big)=e^{\tau\Lin}\Riesz f$.
Since $\Riesz$ is bounded on $X^k$, it commutes with the integral, and
$(1-\Riesz)\wlin_0(u)=e^{u\Lin}(1-\Riesz)w_0+\int_0^ue^{(u-v)\Lin}(1-\Riesz)F(v)\dd v$.
By Proposition~I.\ref{I-prop:semigroup} and the bound
$\|h\cdot\nabla\Phi\|_{H^k}\le\Cnab(k)\|h\|_{C^k(N')}$ of Lemma~I.\ref{I-lem:uniform}, its
norm in $H^k$ is at most
$C_S(k,\sigma')\big(e^{\sigma'u}A+\Cnab(k)\int_0^ue^{\sigma'(u-v)}(A\jap{v}^p+\eta_b\kappa e^{\sigma v})\dd v\big)$.
Finally
$\int_0^ue^{\sigma'(u-v)}\jap{v}^p\dd v\le\jap{u}^pe^{\sigma'u}/\sigma'$ and
$\int_0^ue^{\sigma'(u-v)}e^{\sigma v}\dd v\le e^{\sigma u}/(\sigma-\sigma')\le e^{\sigma u}/(\sigb-\sigma')$.
\end{proof}

For $x\in(0,1]$ we write $\varsigma(x)=1+\sigb^{-1}\log(1/x)$. Since
$\sigma(G)\ge\sigb$, this number bounds $1+u$ for $u\le\sigma(G)^{-1}\log(1/x)$,
the time in which a coefficient of size $x$ that grows at the rate $\sigma(G)$
reaches size $1$. For an integer $s\ge3$
and a number $\sigma'\in(\frac9{200},\sigb)$ we put
\begin{equation}
\label{eq:W-const}
\begin{gathered}
p=3,\qquad N_G=\min\big\{n\in\Z:n\ge2,\ n\ge\Lam(s)/\sigb+1\big\},\qquad
M=s+N_G+2,\\ N_W=(p+2)N_G+1,\qquad
\bar C_S=\max_{0\le k\le M}C_S(k,\sigma'),\qquad
C_{\mathrm{lin}}=\max_{2\le k\le M}C_{\mathrm{lin}}(k),\\
C_{\mathrm{bil}}=\max_{2\le k\le M}\max\big(C_B(k),C_h(k),\Cnab(k)\big),\\
C_\phi=\max_{G\in[0,\Gone]}\|\phist(G)\|_{H^M},\qquad
\bar\rho=\sup_{u\ge0}\jap{u}^pe^{-(\sigb-\sigma')u},
\end{gathered}
\end{equation}
with $C_S$ of Proposition~I.\ref{I-prop:semigroup}, $C_B$ and $C_h$ of
Lemma~\ref{lem:bilinear}, $\Cnab$ of Lemma~I.\ref{I-lem:uniform} and
$C_{\mathrm{lin}}(k)$ of Lemma~\ref{lem:linear-part}. The maximum defining $C_\phi$
exists and is finite by Lemma~I.\ref{I-lem:uniform}, and $\bar\rho<\infty$ because
$\sigb>\sigma'$. All these numbers depend only on $(\Phi,\Gone,s,\sigma')$. As in
Theorem~\ref{thm:quiet}, the order $N_G$ of the approximate solution is chosen
large, so that its residual, which is one order higher, is small enough on the
window of Lemma~\ref{lem:window} to absorb the growth $e^{\Lam(s)u}$ allowed by
Lemma~\ref{lem:energy} (Equation~\eqref{eq:W-exp} below).

Lemma~\ref{lem:window} is the main step in the proof of Theorem~\ref{thm:active}.
It treats an active ring whose initial deviation and running coefficient are
bounded by a fixed multiple $C_\natural$ of a number $\kappa_1$, and whose field is
bounded accordingly. It does so on the window of time in which a component of
initial size $\kappa_1$, growing at the rate $\sigma(G)$, stays below $\delta_0$.

\begin{lemma}
\label{lem:window}
Assume \textup{(N1)}--\textup{(N3)}. Let $s\ge3$ be an integer and
$\sigma'\in(\frac9{200},\sigb)$, let $p$, $N_G$, $M$ and $N_W$ be as in
Equation~\eqref{eq:W-const}, and let $C_\natural\ge1$. There are $\delta_W\in(0,\frac12]$, $\kappa_W\in(0,1]$ and
$C_W\ge1$, depending only on $(\Phi,\Gone,s,\sigma',C_\natural)$, such that the
following holds. Let $G\in[0,\Gone]$, $\sigma=\sigma(G)$, $\lambda=\lamst(G)$,
$\kappa_1\in(0,\kappa_W]$, $\delta_0\in(0,\delta_W]$, $\eta_b\in[0,1]$,
$\hat T_1=\sigma^{-1}\log(\delta_0/\kappa_1)$ and $\bar u\in(0,\hat T_1]$. Let
$\omega$, $\omega_{\mathrm r},\omega_{\mathrm i},\omega_{\mathrm o}$ and $h$ be as in
Lemma~\ref{lem:reduction} on $[0,\bar u]$ with this $G$, put
$w=\omega_{\mathrm r}-\Phi$, and let $c$ be the running coefficient of Equation~\eqref{eq:running}
with $w_0=w(0)$. Assume that $\|w(0)\|_{H^{M-2}}\le C_\natural\kappa_1$, that
$\supp(\Phi+w(0))\subset\{(\frac78-\frac1{64})a\le|y|\le\frac54+\frac a{64}\}$, and
that for $u\in[0,\bar u]$,
\[
\|h(u)\|_{C^M(N')}\le C_\natural\kappa_1\jap{u}^p+\eta_b\kappa_1e^{\sigma u},\qquad
|c(u)|\le C_\natural\kappa_1 .
\]
Put $\hat a_1(u)=\kappa_1e^{\sigma u}$ and $\rho_1(u)=\jap{u}^pe^{-(\sigma-\sigma')u}$.
Then for every $u\in[0,\bar u]$, every trajectory of $\omega_{\mathrm r}$ moves
radially by at most $C_W\delta_0+C_W\kappa_1\varsigma(\kappa_1)^{N_W}\le a/32$ on
$[0,u]$, and
\[
\big\|w(u)-2\Re\big(c(u)e^{\lambda u}\phist\big)\big\|_{H^s}\le C_W\hat a_1(u)\big[\rho_1(u)+\hat a_1(u)+\kappa_1\varsigma(\kappa_1)^{N_W}+\eta_b\big].
\]
\end{lemma}
\begin{proof}
Write $\Lam=\Lam(s)$ and $\Lin=\Lin(G)$. Since $N_G\ge\Lam/\sigb+1$ and
$\Lam\ge2\sigma\ge2\sigb$ (Lemma~\ref{lem:energy}),
\begin{equation}
\label{eq:W-exp}
(N_G+1)\sigma-\Lam\ge2\sigma,\qquad\Lam-\sigma\ge\sigma\ge\sigb,\qquad N_G+1-\Lam/\sigma\ge2 .
\end{equation}

\emph{Constants.} With the numbers of Equation~\eqref{eq:W-const}, put
$C^{(1)}=2C_\natural C_\phi+C_{\mathrm{lin}}(C_\natural\bar\rho+1)$ and, for
$2\le j\le N_G$,
\[
C^{(j)}=2^{j-1}\max(1,\sigb^{-1})\,\bar C_SC_{\mathrm{bil}}\Big(\sum_{i+l=j}C^{(i)}C^{(l)}+C^{(j-1)}\Big).
\]
Put $C_{\mathrm{res}}=C_{\mathrm{bil}}\big(\sum_{i,l\le N_G,\,i+l>N_G}C^{(i)}C^{(l)}+2C^{(N_G)}\big)$,
$C_{\mathrm{en}}=2^{N_G-1}C_{\mathrm{res}}/\sigb$, $S_W=\sum_{j=1}^{N_G}C^{(j)}+C_{\mathrm{en}}$,
\[
C_W=2\big(1+\sigb^{-1}\big)\big(1+C_{\mathrm{BS}}\big)C_\natural^{N_G}\big(C_{\mathrm{lin}}C_\natural+S_W\big),\qquad
\delta_W=\min\Big(\frac12,\frac1{S_W},\frac a{64C_W}\Big),
\]
and let $\kappa_W$ be the largest number in $(0,1]$ such that
$C_Wx\varsigma(x)^{N_W}\le a/64$ for every $x\in(0,\kappa_W]$. It exists by the
argument given for $\kappa_0(C_0)$ in the proof of Theorem~\ref{thm:active}, because
$x\varsigma(x)^{N_W}\to0$ as $x\to0^+$ and the inequality is non-strict between
continuous functions. Taking $t=0$ in
Proposition~I.\ref{I-prop:semigroup} gives $C_S\ge1$, so $C_{\mathrm{lin}}\ge1$ and
$C_W\ge C_\natural$.

Now fix the data, and put $\varepsilon_W=C_\natural\kappa_1$,
$\mathfrak b(u)=\hat a_1(u)+\varepsilon_W\jap{u}^{p+2}$ and $k_j=M-1-j$ for
$1\le j\le N_G$, so $k_1=M-2$, $k_{N_G}=s+1$, $k_j+1=k_{j-1}$. The function
$\mathfrak b$ plays the role of $\maj$ in the proof of Theorem~\ref{thm:quiet}: the
terms $w^{[j]}$ below are bounded by multiples of $\hat a_1\mathfrak b^{j-1}$. For
$u\le\bar u\le\hat T_1\le\sigma^{-1}\log(1/\kappa_1)$ we have $\jap{u}\le1+u\le\varsigma(\kappa_1)$,
$\hat a_1(u)\le\delta_0$ and $\kappa_1\le\hat a_1(u)$. By the choice of $\kappa_W$ and
$\delta_W$, for $u\in[0,\bar u]$,
\begin{equation}
\label{eq:W-small}
\mathfrak b(u)\le\delta_0+C_\natural\kappa_1\varsigma(\kappa_1)^{p+2}\le1,\qquad
\|h(u)\|_{C^{s+1}(N')}\le C_\natural\kappa_1\varsigma(\kappa_1)^p+\delta_0\le1 .
\end{equation}

\emph{Approximate solution.} Let $w^{[j]}$ be the functions of
Equation~\eqref{eq:hier} with $J=N_G$ and $w^{[1]}=\wlin_0$, the function of
Lemma~\ref{lem:running} with $w_0=w(0)$. This choice of $w^{[1]}$ is admissible in
Equation~\eqref{eq:hier} because $\supp(\Phi+w(0))\subset N$, and the functions
$w^{[j]}$ are regular by Lemmas~\ref{lem:duhamel} and \ref{lem:bilinear}. We
claim that
\begin{equation}
\label{eq:W-hier}
\|w^{[j]}(u)\|_{H^{k_j}}\le C^{(j)}\hat a_1(u)\,\mathfrak b(u)^{j-1}\qquad(u\in[0,\bar u],\ 1\le j\le N_G).
\end{equation}
For $j=1$, Lemma~\ref{lem:linear-part} with $k=M-2$, $A=C_\natural\kappa_1$ and
$\kappa=\kappa_1$ gives
$\|(1-\Riesz)w^{[1]}(u)\|_{H^{M-2}}\le C_{\mathrm{lin}}\hat a_1(u)\big(C_\natural\rho_1(u)+\eta_b\big)$,
since $C_\natural\kappa_1\jap{u}^pe^{\sigma'u}=C_\natural\hat a_1(u)\rho_1(u)$; and
$\|2\Re(c(u)e^{\lambda u}\phist)\|_{H^{M-2}}\le2C_\phi C_\natural\hat a_1(u)$. As
$\rho_1\le\bar\rho$ and $\eta_b\le1$, this is Equation~\eqref{eq:W-hier} for $j=1$.
Let $j\ge2$. By Lemma~\ref{lem:bilinear}, Equation~\eqref{eq:W-hier} for the
indices less than $j$, and
$\|h(v)\|_{C^{k_j}(N')}\le\varepsilon_W\jap{v}^p+\hat a_1(v)$,
\[
\|g^{[j]}(v)\|_{H^{k_j}}\le C_{\mathrm{bil}}\Big(\sum_{i+l=j}C^{(i)}C^{(l)}+C^{(j-1)}\Big)\hat a_1(v)\big(\hat a_1(v)+\varepsilon_W\jap{v}^p\big)\mathfrak b(v)^{j-2}.
\]
Since $e^{\sigma(u-v)}\hat a_1(v)=\hat a_1(u)$ and
$\|e^{(u-v)\Lin}\|\le\bar C_Se^{\sigma(u-v)}$, it suffices to show that
\[
\int_0^u\big(\hat a_1+\varepsilon_W\jap{v}^p\big)\mathfrak b^{j-2}\dd v\le2^{j-1}\max(1,\sigb^{-1})\,\mathfrak b(u)^{j-1}.
\]
Expand $\mathfrak b^{j-2}=\sum_{i+q=j-2}\binom{j-2}i\hat a_1^i(\varepsilon_W\jap{v}^{p+2})^q$,
and use that $\hat a_1$ and $\jap{\cdot}$ are nondecreasing. Then
\[
\begin{gathered}
\int_0^u\hat a_1^{i+1}(\varepsilon_W\jap{v}^{p+2})^q\dd v\le\frac{\hat a_1(u)^{i+1}(\varepsilon_W\jap{u}^{p+2})^q}{(i+1)\sigma},\\
\int_0^u\varepsilon_W\jap{v}^p\hat a_1^i(\varepsilon_W\jap{v}^{p+2})^q\dd v\le\frac{\hat a_1(u)^i(\varepsilon_W\jap{u}^{p+2})^{q+1}}{i\sigma}\quad(i\ge1),\\
\int_0^u\varepsilon_W^{j-1}\jap{v}^{p+(p+2)(j-2)}\dd v\le(\varepsilon_W\jap{u}^{p+2})^{j-1},
\end{gathered}
\]
where the second line uses $\jap{u}^p\le\jap{u}^{p+2}$ and the third uses
$u\le\jap{u}$. Each right side is a monomial of degree $j-1$ in
$\hat a_1(u)$ and $\varepsilon_W\jap{u}^{p+2}$, hence at most
$\max(1,\sigb^{-1})\mathfrak b(u)^{j-1}$, and there are $2\cdot2^{j-2}$ terms counted with
the binomial weights. This proves Equation~\eqref{eq:W-hier}.

The residual is given by Equation~\eqref{eq:hier-res} with $J=N_G$. Since
$\mathfrak b\le1$ and $\hat a_1\le\mathfrak b$, we have
$\hat a_1^2\mathfrak b^{i+l-2}\le\hat a_1\mathfrak b^{N_G}$ for $i+l>N_G$, and
$\|h\|_{C^s(N')}\le\varepsilon_W\jap{u}^p+\hat a_1\le2\mathfrak b$. So
\begin{equation}
\label{eq:W-res}
\|R_{\mathrm{app}}(u)\|_{H^s}\le C_{\mathrm{res}}\hat a_1(u)\mathfrak b(u)^{N_G},\qquad
\|w_{\mathrm{app}}(u)\|_{H^{s+1}}\le\Big(\sum_jC^{(j)}\Big)\hat a_1(u)\le S_W\delta_0\le1 .
\end{equation}

\emph{Continuity argument.} Let $u_*$ be the largest $u\in[0,\bar u]$ with
$\delta_{\mathrm r}(u)\le a/16$ and $\|w(v)\|_{H^s}\le1$ for $v\in[0,u]$; at $u=0$
both hold, since $\delta_{\mathrm r}(0)=0$ and, as $s\le M-2$,
$\|w(0)\|_{H^s}\le C_\natural\kappa_1\le a/64$ (recall $C_W\ge C_\natural$, $\varsigma\ge1$
and the choice of $\kappa_W$). For $u\le u_*$ the
support of $\omega_{\mathrm r}(u)$ lies in
$\{\frac{51}{64}a\le|y|\le\frac54+\frac5{64}a\}\subset\Nflat\subset N$
(Lemma~\ref{lem:annuli}). So
Lemma~\ref{lem:reduction}\textup{(3)} and Lemma~\ref{lem:energy} apply on $[0,u_*]$,
by Equations~\eqref{eq:W-small} and \eqref{eq:W-res}. As $v=w-w_{\mathrm{app}}$
vanishes at $u=0$, and
$\mathfrak b^{N_G}\le2^{N_G-1}\big(\hat a_1^{N_G}+(\varepsilon_W\jap{\cdot}^{p+2})^{N_G}\big)$,
\[
\begin{aligned}
\|v(u)\|_{H^s}\le2^{N_G-1}C_{\mathrm{res}}\Big(&\int_0^ue^{\Lam(u-v)}\hat a_1(v)^{N_G+1}\dd v\\
&+(\varepsilon_W\jap{u}^{p+2})^{N_G}\int_0^ue^{\Lam(u-v)}\hat a_1(v)\dd v\Big).
\end{aligned}
\]
By Equation~\eqref{eq:W-exp}, the first integral is at most
$\hat a_1(u)^{N_G+1}/((N_G+1)\sigma-\Lam)\le\hat a_1(u)^{N_G+1}/(2\sigb)$, and the second is
at most $\kappa_1e^{\Lam u}/(\Lam-\sigma)\le\kappa_1e^{\Lam u}/\sigb$. For $u\le\hat T_1$,
$\kappa_1e^{\Lam u}\le\kappa_1^{1-\Lam/\sigma}\delta_0^{\Lam/\sigma}$. Hence, for $u\le u_*$,
\begin{equation}
\label{eq:W-v}
\|v(u)\|_{H^s}\le C_{\mathrm{en}}\big(\hat a_1(u)^{N_G+1}+\Theta_W(u)\big),\qquad
\Theta_W(u)=(\varepsilon_W\jap{u}^{p+2})^{N_G}\kappa_1^{1-\Lam/\sigma}\delta_0^{\Lam/\sigma},
\end{equation}
and, by Equation~\eqref{eq:W-exp}, $\kappa_1\le1$ and $\delta_0\le1$,
\begin{equation}
\label{eq:W-theta}
\Theta_W(u)\le C_\natural^{N_G}\varsigma(\kappa_1)^{(p+2)N_G}\kappa_1^2\le C_\natural^{N_G}\kappa_1\varsigma(\kappa_1)^{N_W}\cdot\kappa_1 .
\end{equation}

We collect the bounds for $u\le u_*$. Write
$w-2\Re(ce^{\lambda u}\phist)=(1-\Riesz)w^{[1]}+\sum_{j\ge2}w^{[j]}+v$. By the bound on
$(1-\Riesz)w^{[1]}$ above, Equation~\eqref{eq:W-hier} with $\mathfrak b\le1$,
Equations~\eqref{eq:W-v} and \eqref{eq:W-theta}, $\hat a_1^{N_G+1}\le\hat a_1^2$ and $\Theta_W\le\hat a_1\Theta_W/\kappa_1$,
\[
\begin{aligned}
\big\|w-2\Re(ce^{\lambda u}\phist)\big\|_{H^s}\le{}&C_{\mathrm{lin}}\hat a_1(C_\natural\rho_1+\eta_b)\\
&+S_W\hat a_1\big(2\hat a_1+C_\natural\kappa_1\varsigma(\kappa_1)^{p+2}+C_\natural^{N_G}\kappa_1\varsigma(\kappa_1)^{N_W}\big),
\end{aligned}
\]
which is at most $C_W\hat a_1\big[\rho_1+\hat a_1+\kappa_1\varsigma(\kappa_1)^{N_W}+\eta_b\big]$.
Similarly,
$\|w(u)\|_{H^s}\le S_W\big(\hat a_1(u)+\Theta_W(u)\big)\le C_W\big(\delta_0+\kappa_1\varsigma(\kappa_1)^{N_W}\big)\le\tfrac a{32}$.
For the displacement, Lemma~\ref{lem:reduction}\textup{(2)} and
Lemma~\ref{lem:bilinear}\textup{(4)} give, for $u\le u_*$,
\[
\begin{aligned}
\delta_{\mathrm r}(u)&\le\int_0^u\big(\|h\|_{C^0(N')}+C_{\mathrm{BS}}\|w\|_{H^s}\big)\\
&\le C_\natural\kappa_1\varsigma(\kappa_1)^{p+1}+\frac{\delta_0}{\sigb}+C_{\mathrm{BS}}\Big(\frac{S_W\delta_0}{\sigb}+C_{\mathrm{en}}C_\natural^{N_G}\varsigma(\kappa_1)^{N_W}\kappa_1^2\Big),
\end{aligned}
\]
where we used $\|w\|_{H^s}\le\sum_j\|w^{[j]}\|_{H^s}+\|v\|_{H^s}\le S_W\hat a_1+C_{\mathrm{en}}\Theta_W$,
$\int_0^u\hat a_1\le\hat a_1(u)/\sigma\le\delta_0/\sigb$, and
$\int_0^u\Theta_W\le\bar u\,C_\natural^{N_G}\varsigma(\kappa_1)^{(p+2)N_G}\kappa_1^2$ with
$\bar u\le\varsigma(\kappa_1)$. The right side is at most
$C_W\delta_0+C_W\kappa_1\varsigma(\kappa_1)^{N_W}\le a/32$. If $u_*<\bar u$, continuity
would give a larger $u$ with $\delta_{\mathrm r}(u)\le a/16$ and $\|w\|_{H^s}\le1$ on $[0,u]$, a
contradiction. So $u_*=\bar u$, and the bounds just proved hold on $[0,\bar u]$.
\end{proof}

\begin{theorem}
\label{thm:active}
Assume \textup{(N1)}--\textup{(N3)}. Fix an integer $s\ge3$, a number
$\sigma'\in(\frac9{200},\sigb)$, and put $p=3$ $(=m)$. Let $N_G$ be the least integer
with $N_G\ge2$ and $N_G\ge\Lam(s)/\sigb+1$, put $M=s+N_G+2$ (so $M\ge s+4$), and
put $\etmax=1$. There are $\dH\in(0,1)$, $C_2\ge1$ and $C_3>0$, depending only on
$(\Phi,\Gone,s,\sigma')$, and for every $C_0\ge1$ a number $\kappa_0(C_0)>0$,
depending only on $C_0$ and $(\Phi,\Gone,s,\sigma')$, such that the following holds.

Let $G\in[0,\Gone]$, $\sigma=\sigma(G)$, $\lambda=\lamst(G)$, $\kappa>0$,
$\eta_b\in[0,\etmax]$, $\delta_0\in(0,\dH]$, $\hat a(u)=\kappa e^{\sigma u}$ and
$\hat T=\sigma^{-1}\log(\delta_0/\kappa)$. Let $t_1=t_1(C_0,G)$ be the least
$t\ge0$ such that $2^{p/2}C_0\jap{t'}^pe^{-(\sigma-\sigma')t'}\le1$ for all $t'\ge t$,
and define the \textbf{transient}
\[
\rhohat(u)=C_0\jap{u}^pe^{-(\sigma-\sigma')u}\ (0\le u<t_1),\qquad
\rhohat(u)=\jap{u-t_1}^pe^{-(\sigma-\sigma')(u-t_1)}\ (u\ge t_1).
\]
Let $\bar u>0$, let $\omega$, $\omega_{\mathrm r},\omega_{\mathrm i},\omega_{\mathrm o}$
and $h$ be as in Lemma~\ref{lem:reduction} on $[0,\bar u]$ with this $G$, put
$w(u)=\omega_{\mathrm r}(u)-\Phi$, and let $c$ be the running coefficient of
Equation~\eqref{eq:running} with $w_0=w(0)$. Assume:
\begin{description}[font=\normalfont]
\item[\textup{(Init)}] $\|w(0)\|_{H^{M-1}}\le C_0\kappa$ and $\supp(\Phi+w(0))\subset\NNm$;
\item[\textup{(Field)}] $h=h_1+h_2$ on $[0,\bar u]$ with
  $\|h_1(u)\|_{C^M(N')}\le C_0\kappa\jap{u}^p$ and
  $\|h_2(u)\|_{C^M(N')}\le\eta_b\,\hat a(u)$;
\item[\textup{(Coef)}] $|c(u)|\le\kappa\big(2+C_0\jap{u}^pe^{-\sigma u}\big)$ on $[0,\bar u]$;
\item[\textup{(Small)}] $\kappa\le\kappa_0(C_0)$.
\end{description}
Then for every $u\in[0,\min(\bar u,\hat T)]$:
\begin{enumerate}[label=\textup{(\alph*)}]
\item $\supp(\Phi+w(u))\subset\NN$;
\item $\big\|w(u)-2\Re\big(c(u)e^{\lambda u}\phist\big)\big\|_{H^s}\le
  C_2\hat a(u)\big[\rhohat(u)+\hat a(u)+\kappa^{1/2}+\eta_b\big]$ and
  $\|w(u)\|_{H^s}\le C_2\hat a(u)(1+\rhohat(u))$.
\end{enumerate}
Moreover, let $\sigma_{\max}=\max_{G\in[0,\Gone]}\sigma(G)$, and let $\bar t_1(C_0)$ be the
least $t\ge0$ such that
\[
2^{p/2}C_0\jap{t'}^pe^{-(\sigb-\sigma')t'}\le1\qquad\text{for all }t'\ge t .
\]
Then $t_1(C_0,G)\le\bar t_1(C_0)$, and for all $u\ge0$,
\[
\hat a(u)\rhohat(u)\le C(C_0)\,\kappa\jap{u}^pe^{\sigma'u},\qquad
C(C_0)=\max\big(C_0,e^{\sigma_{\max}\bar t_1(C_0)}\big);
\]
the number $C(C_0)$ depends only on $C_0$, $\sigma'$, $p$, $\sigb$ and $\sigma_{\max}$.
\end{theorem}
\begin{proof}
Throughout,
$\Lin=\Lin(G)$, $F=-h\cdot\nabla\Phi$, and $\wlin$ is the function $\wlin_0$ of
Lemma~\ref{lem:running} with $w_0=w(0)$; it is an admissible choice of $w^{[1]}$
in Equation~\eqref{eq:hier}, because $\supp(\Phi+w(0))\subset\NNm\subset N$ by
(Init). The proof has two stages. On
the initial layer $[0,\min(\bar u,t_1)]$, whose length is at most $\bar t_1(C_0)$,
we compare $w$ with $\wlin$ by Lemma~\ref{lem:energy}. At the time $t_1$ the
transient has decayed, and we restart the problem there and apply
Lemma~\ref{lem:window} on the rest of the interval.

\emph{Constants.} Let $C_\phi$, $\bar C_S$, $C_{\mathrm{bil}}$ and
$C_{\mathrm{lin}}$ be the numbers of Equation~\eqref{eq:W-const}, $\|\ellst\|$ and
$\Cp$ as in Lemma~I.\ref{I-lem:uniform}, and put
$C_\natural=\max(4,\,6C_\phi+2C_{\mathrm{lin}}+1)$; this choice makes the restarted
problem below satisfy the assumptions of Lemma~\ref{lem:window}. Let $\delta_W,\kappa_W,C_W$ be the
constants of Lemma~\ref{lem:window} for this $C_\natural$, and put
\[
C_2=6C_\phi+3\big(C_W+C_{\mathrm{lin}}+1\big),\qquad C_3=(1+\Cp)C_2,\qquad
\dH=\min\Big(\delta_W,\frac a{32C_2}\Big).
\]
These depend only on $(\Phi,\Gone,s,\sigma')$. Now let $C_0\ge1$, write
$\bar t_1=\bar t_1(C_0)$, $\Lambda'=\Lam(M-1)$, $\Lambda''=\Lam(M-2)$,
$C_{\mathrm f}=C_0\jap{\bar t_1}^p+e^{\sigma_{\max}\bar t_1}$, and define the
constants of the initial layer by
\[
\begin{gathered}
C^{\mathrm{lay}}_1=e^{\Lambda'\bar t_1}\big(C_0+C_{\mathrm{bil}}C_{\mathrm f}\bar t_1\big),\qquad
C^{\mathrm{lay}}_2=\bar t_1\big(C_{\mathrm f}+C_{\mathrm{BS}}C^{\mathrm{lay}}_1\big),\\
C^{\mathrm{lay}}_3=\bar C_Se^{\sigma_{\max}\bar t_1}\big(C_0+C_{\mathrm{bil}}C_{\mathrm f}\bar t_1\big),\qquad
C^{\mathrm{lay}}_4=\bar t_1e^{\Lambda''\bar t_1}C_{\mathrm{bil}}\big(C^{\mathrm{lay}}_3+C_{\mathrm f}\big)C^{\mathrm{lay}}_3 .
\end{gathered}
\]
The smallness of $\kappa$ is imposed by three groups of conditions: (L1) is used on
the initial layer, (L2) to bound the terms that contain the error $v_1=w-\wlin$ of
the initial layer, and (L3) to apply Lemma~\ref{lem:window} after the restart.
Let $\kappa_0(C_0)$ be the largest $x\in(0,(a/32)^2]$ such that every
$\kappa\in(0,x]$ satisfies:
\begin{enumerate}[label=\textup{(L\arabic*)}]
\item $C_{\mathrm f}\kappa\le1$, $C^{\mathrm{lay}}_1\kappa\le\frac12$, $C^{\mathrm{lay}}_3\kappa\le1$ and
  $C^{\mathrm{lay}}_2\kappa\le\min(\frac a{128},\frac12\kappa^{1/2})$;
\item $(1+\|\ellst\|)C^{\mathrm{lay}}_4\kappa\le1$ and
  $4(1+C_\phi\|\ellst\|)C^{\mathrm{lay}}_4\kappa^{1/2}\le1$;
\item $e^{\sigma_{\max}\bar t_1}\kappa\le\kappa_W$ and
  $C_We^{\sigma_{\max}\bar t_1}\kappa\,\varsigma(\kappa)^{N_W}\le\frac12\kappa^{1/2}$.
\end{enumerate}
This number exists. Each inequality in (L1)--(L3) is non-strict between continuous
functions of $\kappa\in(0,1]$, so the set $S$ of $\kappa\in(0,(a/32)^2]$ satisfying all of
them is closed in $(0,(a/32)^2]$; and $S\supset(0,\epsilon]$ for some $\epsilon>0$,
because in each inequality the ratio of the left side to the right side tends to $0$
as $\kappa\to0^+$. So the supremum $x_0$ of the $x$ with $(0,x]\subset S$ is positive,
$(0,x_0)\subset S$, and $x_0\in S$ by closedness. The number $\kappa_0(C_0)$ depends
only on $C_0$ and $(\Phi,\Gone,s,\sigma')$, and by (Small) the $\kappa$ of the theorem
satisfies (L1)--(L3) and $\kappa\le(a/32)^2$.

\emph{The time $t_1$.} Since $\sigma\ge\sigb$,
$2^{p/2}C_0\jap{t'}^pe^{-(\sigma-\sigma')t'}\le2^{p/2}C_0\jap{t'}^pe^{-(\sigb-\sigma')t'}$, so
$t_1\le\bar t_1$. For $u\ge t_1$, the definition of $t_1$ with $t'=u$ gives
\begin{equation}
\label{eq:A-t1}
C_0\jap{u}^pe^{-\sigma u}\le2^{-p/2}e^{-\sigma'u}\le1,\qquad\text{hence }|c(u)|\le3\kappa\text{ by (Coef)}.
\end{equation}
For $u<t_1$,
$\hat a(u)\rhohat(u)=C_0\kappa\jap{u}^pe^{\sigma'u}$. For $u\ge t_1$,
$\hat a(u)\rhohat(u)=\kappa e^{\sigma t_1}\jap{u-t_1}^pe^{\sigma'(u-t_1)}\le e^{\sigma_{\max}\bar t_1}\kappa\jap{u}^pe^{\sigma'u}$.
This proves the last assertion of the theorem.

\emph{The initial layer $[0,\bar u_1]$, $\bar u_1=\min(\bar u,t_1)$.} Since
$t_1\le\bar t_1$ and $\eta_b\le1$, (Field) gives
$\|h(u)\|_{C^M(N')}\le C_{\mathrm f}\kappa\le1$ there. Let $u_*$ be the largest
$u\in[0,\bar u_1]$ such that $\delta_{\mathrm r}(u)\le a/64$ and $\|w(v)\|_{H^{M-1}}\le1$ for
$v\in[0,u]$; both hold at $u=0$ by (Init) and (L1). For $u\le u_*$ the support of
$\omega_{\mathrm r}(u)$ lies in $\{(\frac78-\frac1{64})a\le|y|\le\frac54+\frac a{64}\}\subset N$
(Lemma~\ref{lem:annuli}), so Lemma~\ref{lem:reduction}\textup{(3)} applies on $[0,u_*]$. Lemma~\ref{lem:energy}
at level $M-1$ with $w_{\mathrm{app}}=0$, for which $R_{\mathrm{app}}=h\cdot\nabla\Phi$, and
Lemma~\ref{lem:bilinear}\textup{(3)} give
\[
\|w(u)\|_{H^{M-1}}\le e^{\Lambda'u}C_0\kappa+\int_0^ue^{\Lambda'(u-v)}C_{\mathrm{bil}}C_{\mathrm f}\kappa\dd v\le C^{\mathrm{lay}}_1\kappa\le\tfrac12 ,
\]
and Lemma~\ref{lem:reduction}\textup{(2)} and Lemma~\ref{lem:bilinear}\textup{(4)} give
$\delta_{\mathrm r}(u)\le\int_0^u\big(\|h\|_{C^0(N')}+C_{\mathrm{BS}}\|w\|_{H^1}\big)\le C^{\mathrm{lay}}_2\kappa\le a/128$.
If $u_*<\bar u_1$, continuity would give a larger admissible $u$; so $u_*=\bar u_1$.
Next, $\wlin\in C([0,\bar u];X^{M-1})\cap C^1([0,\bar u];X^{M-2})$ is real with
$\partial_u\wlin=\Lin\wlin+F$ (Lemma~\ref{lem:duhamel}), and
Proposition~I.\ref{I-prop:semigroup} with Lemma~\ref{lem:bilinear} gives
$\|\wlin(u)\|_{H^{M-1}}\le C^{\mathrm{lay}}_3\kappa\le1$ for $u\le\bar u_1$. So
Lemma~\ref{lem:energy} at level $M-2$ applies to $w_{\mathrm{app}}=\wlin$ on
$[0,\bar u_1]$. Its residual, Equation~\eqref{eq:hier-res} with $J=1$, is
$R_{\mathrm{app}}=-B(\wlin,\wlin)+h\cdot\nabla\wlin$, and
$\|R_{\mathrm{app}}\|_{H^{M-2}}\le C_{\mathrm{bil}}(C^{\mathrm{lay}}_3+C_{\mathrm f})C^{\mathrm{lay}}_3\kappa^2$.
Since $v_1:=w-\wlin$ vanishes at $u=0$,
\begin{equation}
\label{eq:A-layer}
\|v_1(u)\|_{H^{M-2}}\le C^{\mathrm{lay}}_4\kappa^2,\qquad \delta_{\mathrm r}(u)\le C^{\mathrm{lay}}_2\kappa\qquad(0\le u\le\bar u_1).
\end{equation}

\emph{Proof of (b) for $u\le\min(\bar u_1,\hat T)$.} By
Lemma~\ref{lem:linear-part} with $k=s$ and $A=C_0\kappa$, and since
$C_0\kappa\jap{u}^pe^{\sigma'u}\le\hat a(u)\rhohat(u)$ for $u\le t_1$ (with
equality for $u<t_1$, while at $u=t_1$ the left side is at most
$2^{-p/2}\hat a(t_1)$ by the definition of $t_1$ and $\rhohat(t_1)=1$),
\[
\begin{gathered}
\wlin(u)=2\Re\big(c(u)e^{\lambda u}\phist\big)+(1-\Riesz)\wlin(u),\\
\|(1-\Riesz)\wlin(u)\|_{H^s}\le C_{\mathrm{lin}}\hat a(u)\big(\rhohat(u)+\eta_b\big).
\end{gathered}
\]
By Equation~\eqref{eq:A-layer} and (L2),
$\|v_1(u)\|_{H^s}\le C^{\mathrm{lay}}_4\kappa^2\le\kappa^{1/2}\kappa\le\kappa^{1/2}\hat a(u)$. Hence
$\|w-2\Re(ce^{\lambda u}\phist)\|_{H^s}\le(C_{\mathrm{lin}}+1)\hat a\big[\rhohat+\kappa^{1/2}+\eta_b\big]$.
By (Coef), $|c(u)|e^{\sigma u}\le2\hat a(u)+C_0\kappa\jap{u}^p\le\hat a(u)(2+\rhohat(u))$,
so $\|w(u)\|_{H^s}\le2C_\phi(2+\rhohat)\hat a+(C_{\mathrm{lin}}+1)(\rhohat+2)\hat a\le C_2\hat a(1+\rhohat)$.
The displacement is at most $C^{\mathrm{lay}}_2\kappa\le\frac12\kappa^{1/2}$ by (L1). This
proves (b), and bounds the displacement by $\frac12\kappa^{1/2}$, when
$\min(\bar u,\hat T)\le t_1$.

\emph{Restart.} Assume now that $t_1<u_e:=\min(\bar u,\hat T)$. The function
$\tilde\omega(u')=\omega(t_1+u')$ is the solution of Theorem~\ref{thm:smooth} with the
$3$-fold rotationally symmetric datum $\omega(t_1)$, by Theorem~\ref{thm:smooth}\textup{(5)}, and its
flow is $\tilde\flow_{u'}=\flow_{t_1+u'}\circ\flow_{t_1}^{-1}$. With
$\tilde\omega_\bullet(0)=\omega_\bullet(t_1)$ we get $\tilde\omega_\bullet(u')=\omega_\bullet(t_1+u')$,
the same $G$ (since $\int\omega_{\mathrm i}(t_1)=2\pi G$), and the field
$\tilde h(u')=h(t_1+u')$. So the restarted configuration satisfies the assumptions of
Lemma~\ref{lem:reduction} on $[0,\bar u-t_1]$, and its perturbation is
$\tilde w(u')=w(t_1+u')$. Its running coefficient is
$\tilde c(u')=\ellst(w(t_1))+\int_0^{u'}e^{-\lambda v}\ellst(F(t_1+v))\dd v$. By
Equation~\eqref{eq:running} and Lemma~\ref{lem:running},
$e^{\lambda t_1}c(t_1+u')=\ellst(\wlin(t_1))+\int_0^{u'}e^{-\lambda v}\ellst(F(t_1+v))\dd v$, so
\begin{equation}
\label{eq:A-restart}
\tilde c(u')=e^{\lambda t_1}c(t_1+u')+\ellst(v_1(t_1)).
\end{equation}

We check the assumptions of Lemma~\ref{lem:window} for the restarted problem,
with $C_\natural$ as above, $\kappa_1=\hat a(t_1)=\kappa e^{\sigma t_1}$, the same
$\delta_0$ and $\eta_b$, and $\bar u'=u_e-t_1$. First,
$\kappa\le\kappa_1\le e^{\sigma_{\max}\bar t_1}\kappa\le\kappa_W$ by (L3),
$\delta_0\le\dH\le\delta_W$, and $\sigma^{-1}\log(\delta_0/\kappa_1)=\hat T-t_1\ge\bar u'$.
By Equation~\eqref{eq:A-t1} at $u=t_1$, $|c(t_1)|e^{\sigma t_1}\le3\kappa_1$ and
$C_0\kappa\jap{t_1}^pe^{\sigma't_1}\le\kappa_1$, and
$\eta_b\kappa e^{\sigma t_1}\le\kappa_1$. So Lemma~\ref{lem:linear-part} at $k=M-2$,
Equation~\eqref{eq:A-layer} and (L2) give
\[
\|w(t_1)\|_{H^{M-2}}\le6C_\phi\kappa_1+C_{\mathrm{lin}}(\kappa_1+\kappa_1)+C^{\mathrm{lay}}_4\kappa^2\le(6C_\phi+2C_{\mathrm{lin}}+1)\kappa_1\le C_\natural\kappa_1 .
\]
Second, by Equation~\eqref{eq:A-layer} the support of $\Phi+w(t_1)$ lies in
$\{(\frac78-\frac1{64})a\le|y|\le\frac54+\frac a{64}\}$. Third, by Peetre's inequality
$\jap{t_1+u'}\le\sqrt2\jap{t_1}\jap{u'}$ and the definition of $t_1$,
\[
C_0\kappa\jap{t_1+u'}^p\le\big[2^{p/2}C_0\jap{t_1}^pe^{-(\sigma-\sigma')t_1}\big]e^{-\sigma't_1}\kappa_1\jap{u'}^p\le\kappa_1\jap{u'}^p,
\]
and $\eta_b\hat a(t_1+u')=\eta_b\kappa_1e^{\sigma u'}$; so
$\|\tilde h(u')\|_{C^M(N')}\le C_\natural\kappa_1\jap{u'}^p+\eta_b\kappa_1e^{\sigma u'}$. Fourth,
by Equations~\eqref{eq:A-restart}, \eqref{eq:A-t1} and \eqref{eq:A-layer} and (L2),
$|\tilde c(u')|\le3e^{\sigma t_1}\kappa+\|\ellst\|C^{\mathrm{lay}}_4\kappa^2\le4\kappa_1\le C_\natural\kappa_1$.

\emph{Proof of (b) for $u\in[t_1,u_e]$.} Lemma~\ref{lem:window} applies; put
$u'=u-t_1$. Then $\kappa_1e^{\sigma u'}=\hat a(u)$ and
$\jap{u'}^pe^{-(\sigma-\sigma')u'}=\rhohat(u)$. Since $\varsigma$ is decreasing and
$\kappa_1\ge\kappa$, (L3) gives $C_W\kappa_1\varsigma(\kappa_1)^{N_W}\le\frac12\kappa^{1/2}$.
By Equation~\eqref{eq:A-restart},
$\tilde c(u')e^{\lambda u'}=c(u)e^{\lambda u}+\ellst(v_1(t_1))e^{\lambda u'}$, and
$|\ellst(v_1(t_1))e^{\lambda u'}|\le\|\ellst\|C^{\mathrm{lay}}_4\kappa^2e^{\sigma u'}\le\|\ellst\|C^{\mathrm{lay}}_4\kappa\,\hat a(u)$.
Since $\|2\Re(X\phist)\|_{H^s}\le2C_\phi|X|$ for $X\in\C$, (L2) gives
\[
\begin{aligned}
\big\|w(u)-2\Re\big(c(u)e^{\lambda u}\phist\big)\big\|_{H^s}
&\le C_W\hat a(u)\big[\rhohat(u)+\hat a(u)+\tfrac12\kappa^{1/2}+\eta_b\big]\\
&\quad+\tfrac12\kappa^{1/2}\hat a(u),
\end{aligned}
\]
which gives the first bound of (b), since $C_W+1\le C_2$. By
Equation~\eqref{eq:A-t1}, the main term $2\Re(c(u)e^{\lambda u}\phist)$ has $H^s$ norm
at most $6C_\phi\hat a(u)$. Since $\hat a+\kappa^{1/2}+\eta_b\le3$, we get
\[
\|w(u)\|_{H^s}\le6C_\phi\hat a+(C_W+1)\hat a(\rhohat+3)\le C_2\hat a(1+\rhohat).
\]
For the displacement, a trajectory of $\omega_{\mathrm r}$ satisfies, with $z=\flow_{t_1}(y)$,
\[
|\flow_u(y)|-|y|=\big(|\tilde\flow_{u'}(z)|-|z|\big)+\big(|\flow_{t_1}(y)|-|y|\big),
\]
so by Lemma~\ref{lem:window} and Equation~\eqref{eq:A-layer} it moves radially by at
most $C_W\delta_0+\frac12\kappa^{1/2}+C^{\mathrm{lay}}_2\kappa\le C_2\delta_0+\kappa^{1/2}$.

\emph{Proof of (a).} In both cases the displacement is at most
$C_2\delta_0+\kappa^{1/2}\le\frac a{32}+\frac a{32}=\frac a{16}$, because
$\delta_0\le\dH\le a/(32C_2)$ and $\kappa\le(a/32)^2$. Since
$\supp\omega_{\mathrm r}(0)\subset\NNm$ by (Init), Lemma~\ref{lem:annuli} gives
$\supp(\Phi+w(u))\subset\nobreak\NN$.
\end{proof}

Recall from the expansion after Definition~I.\ref{I-def:multipole} that the outer
moment $\pout_3$ of a function determines the part in the angular modes $\pm3$ of its
field outside its support. So the outer moment of the deviation of a ring
determines the part in these modes of the field that the ring exerts on the rings
outside it. Recall also from Definition~I.\ref{I-def:eigendata} and the discussion
before it that $\pout_*(G)=\pout_3[\phist(G)]$ is the outer moment of the unstable
eigenfunction.
The next corollary gives this moment for an active ring.

\begin{corollary}
\label{cor:active-moment}
Under the assumptions of Theorem~\ref{thm:active}, put
$q(u)=2\pi\,\pout_3[w(u)]$. Then for every $u\in[0,\min(\bar u,\hat T)]$,
\[
\big|q(u)-2\pi\pout_*(G)\,c(u)e^{\lambda u}\big|\le2\pi C_3\,\hat a(u)\big[\rhohat(u)+\hat a(u)+\kappa^{1/2}+\eta_b\big].
\]
\end{corollary}
\begin{proof}
For $X\in\C$, $\pout_3[2\Re(X\phist)]=X\pout_3[\phist]+\bar X\pout_3[\overline{\phist}]$. The
second term vanishes: $\overline{\phist}$ lies in the angular mode $-3$, and
$\int_0^{2\pi}e^{-3i\theta}e^{-3i\theta}\dd\theta=0$. Since
$\pout_*(G)=\pout_3[\phist(G)]$, linearity gives
$\pout_3[w(u)]-\pout_*(G)c(u)e^{\lambda u}=\pout_3\big[w(u)-2\Re(c(u)e^{\lambda u}\phist)\big]$.
The function in brackets is supported in $N$, so Lemma~I.\ref{I-lem:uniform} and
Theorem~\ref{thm:active}\textup{(b)} bound its modulus by
$\Cp C_2\hat a(u)\big[\rhohat(u)+\hat a(u)+\kappa^{1/2}+\eta_b\big]$. Multiplying by $2\pi$
and using $\Cp C_2\le C_3$ proves the corollary.
\end{proof}

\subsection{The Seeded Ring}

The approximate datum of Definition~\ref{def:approx} replaces its innermost ring by
a rearrangement $\Phi^\kappa$ of $\Phi$ whose deviation from $\Phi$ has an unstable
component of size about $\kappa$. We obtain $\Phi^\kappa$ by composing $\Phi$ with
the time-$\kappa$ flow of a divergence-free field $V$. This flow preserves area, so
$\Phi^\kappa$ is equimeasurable with $\Phi$, and $V$ is chosen with
$V\cdot\nabla\Phi=2\Re\phist(0)$, so that $\Phi^\kappa-\Phi$ is $2\kappa\Re\phist(0)$ to
first order in $\kappa$.

\begin{lemma}
\label{lem:seed}
Assume \textup{(N1)}--\textup{(N3)}, let $s\ge3$ be an integer, and let
$M=s+N_G+2$ be the integer of Theorem~\ref{thm:active} for this $s$. There are a
number $\kL>0$ and functions $\Phi^\kappa$, $\kappa\in(0,\kL]$, all depending only
on $\Phi$, and a number $\Cseed<\infty$, depending only on $\Phi$, $\Gone$ and $s$,
such that for every
$\kappa\in(0,\kL]$: $\Phi^\kappa\in C_c^\infty(\R^2)$, $\Phi^\kappa\ge0$,
$\Phi^\kappa$ is $3$-fold rotationally symmetric and equimeasurable with $\Phi$;
$\supp\Phi^\kappa\subset\NNm$;
$\|\Phi^\kappa-\Phi\|_{H^{M-1}}\le\Cseed\kappa$;
and $|\ellst(0)(\Phi^\kappa-\Phi)|\in[\frac\kappa2,2\kappa]$.
\end{lemma}
\begin{proof}
Let $G=0$, $\lambda=\lamst(0)$, and let $\psi_*$ and $n_*=n_*(0)$ be
as in Theorem~I.\ref{I-thm:evans-simple} and Definition~I.\ref{I-def:eigendata}, for the
ring in ring units and $\ell=3$. By these definitions,
$\phist(0)=n_*^{-1}g_*(r)e^{3i\theta}$ with $g_*=\Mres_{\lambda}\kcoef_3\,\psi_*(\log r)$,
where $\kcoef_3(r)=3i\Phi'(r)/r$ and $\Mres_\lambda=(\lambda+3i\Omega(\cdot;0))^{-1}$
(Definition~I.\ref{I-def:sector-setting}). So the mode-$3$ profile of
$\phist(0)$ is
\[
\vphist(r)=n_*^{-1}\,\frac{3i\,\Phi'(r)}{r}\,\frac{\psi_*(\log r)}{\lambda+3i\Omega(r;0)} .
\]
Put $g_{\mathrm s}(r)=n_*^{-1}\psi_*(\log r)/(\lambda+3i\Omega(r;0))$ for $r>0$. It is
smooth: $\psi_*$ solves a linear equation with smooth coefficients
(Lemma~I.\ref{I-lem:sector}\textup{(3)}), $\Omega(\cdot;0)$ is smooth on $(0,\infty)$, and
$|\lambda+3i\Omega|\ge\Re\lambda>0$. Let $\chiseed\in C_c^\infty(\{a/2<|y|<2\})$ be
radial with $\chiseed=1$ on $\{a\le|y|\le1\}$, and put
$\psi_{\mathrm{seed}}=-2\Re\big(\chiseed(r)\,g_{\mathrm s}(r)\,e^{3i\theta}\big)$ and
$V=\nabla^\perp\psi_{\mathrm{seed}}$. Then $\psi_{\mathrm{seed}}$ is real, smooth,
$3$-fold rotationally symmetric and supported in the compact set
$K_V=\supp\chiseed\subset\{a/2<|y|<2\}$, and $V\in C_c^\infty(\R^2;\R^2)$ is
divergence-free and $3$-fold rotationally symmetric. Let $Y_\kappa=\flow_V(\kappa,0,\cdot)$ be the
time-$\kappa$ flow of $V$, with $\flow_V$ as in Lemma~\ref{lem:flow} for the
time-independent bounded field $V$. By Lemma~\ref{lem:flow}\textup{(a)}, \textup{(b)},
\textup{(d)}, $Y$ is defined for all $\kappa\in\R$ and is $C^\infty$ in $(\kappa,y)$,
each $Y_\kappa$ is a $C^\infty$ diffeomorphism of $\R^2$, and $\det DY_\kappa=1$
because $\divop V=0$. Since $V$ does not depend on time, uniqueness for the flow
gives $Y_{\kappa+\kappa'}=Y_\kappa\circ Y_{\kappa'}$ and, since
$V(\Orot y)=\Orot V(y)$, $Y_\kappa\circ\Orot=\Orot\circ Y_\kappa$. For $y\notin K_V$
the constant path $y$ solves $\dot x=V(x)$, since $V(y)=0$; so $Y_\kappa$ is the identity outside
$K_V$.

Put $\Phi^\kappa=\Phi\circ Y_\kappa$. Then $\Phi^\kappa\in C_c^\infty(\R^2)$ is nonnegative and $3$-fold
rotationally symmetric. It is equimeasurable with $\Phi$: for $c>0$,
$\{\Phi^\kappa>c\}=Y_\kappa^{-1}(\{\Phi>c\})$ has the measure of $\{\Phi>c\}$ because
$Y_\kappa$ preserves area, and $\{\Phi^\kappa<-c\}=\{\Phi<-c\}=\emptyset$. Since
$\Phi$ is radial, $V\cdot\nabla\Phi=(V\cdot\er)\Phi'$, and
$V\cdot\er=-r^{-1}\partial_\theta\psi_{\mathrm{seed}}=2\Re\big(\frac{3i}r\chiseed g_{\mathrm s}e^{3i\theta}\big)$.
On $\supp\Phi'\subset\{a\le|y|\le1\}$ we have $\chiseed=1$, so
$(V\cdot\er)\Phi'=2\Re(\vphist e^{3i\theta})=2\Re\phist(0)$; off $\supp\Phi'$ both
sides vanish, because $\vphist$ is a multiple of $\Phi'$. Thus
\begin{equation}
\label{eq:seed-first}
V\cdot\nabla\Phi=2\Re\phist(0).
\end{equation}

We now expand in $\kappa$. Fix $y$ and put $f_y(t)=\Phi(Y_t(y))$. Then
$f_y'(t)=(V\cdot\nabla\Phi)(Y_t(y))$ and $f_y''(t)=(V\cdot\nabla(V\cdot\nabla\Phi))(Y_t(y))$,
so with $\Psi_{\mathrm s}=V\cdot\nabla(V\cdot\nabla\Phi)\in C_c^\infty(\R^2)$,
\[
\Phi^\kappa-\Phi=\int_0^\kappa(V\cdot\nabla\Phi)\circ Y_t\dd t,\qquad
\Phi^\kappa-\Phi-\kappa V\cdot\nabla\Phi=\int_0^\kappa(\kappa-t)\,\Psi_{\mathrm s}\circ Y_t\dd t .
\]
The derivatives of $Y_t(y)-y$ of every order vanish outside $K_V$ and are
continuous on $[-1,1]\times K_V$, hence bounded there. By the chain rule and
$\det DY_t=1$, for every $k$ there is $C^Y_k$ with $\|g\circ Y_t\|_{H^k}\le C^Y_k\|g\|_{H^k}$
for $g\in C_c^\infty(\R^2)$ and $|t|\le1$. Put $C'_0=\frac12C^Y_0\|\Psi_{\mathrm s}\|_{L^2}$, which depends only on $\Phi$, and
$\Cseed=C^Y_{M-1}\|2\Re\phist(0)\|_{H^{M-1}}$. For $0<\kappa\le1$,
Equation~\eqref{eq:seed-first} and the two identities give
$\|\Phi^\kappa-\Phi-2\kappa\Re\phist(0)\|_{L^2}\le C'_0\kappa^2$ and
$\|\Phi^\kappa-\Phi\|_{H^{M-1}}\le\Cseed\kappa$.

Finally put
$\kL=\min\big(1,\,a/(8(1+\|V\|_{L^\infty})),\,1/(2(1+\|\ellst(0)\|_{(X^0)^*}C'_0))\big)$,
and let $\kappa\in(0,\kL]$. Since $|Y_\kappa^{-1}(x)-x|\le\kappa\|V\|_{L^\infty}\le a/8$ by
Lemma~\ref{lem:flow}\textup{(b)}, \textup{(c)}, and
$\supp\Phi^\kappa=Y_\kappa^{-1}(\supp\Phi)$ with $\supp\Phi\subset\{a\le|y|\le1\}$, we get
$\supp\Phi^\kappa\subset\{\frac78a\le|y|\le1+\frac a8\}\subset\NNm$. The function
$\Phi^\kappa-\Phi$ is supported in $N$ and $3$-fold rotationally symmetric, so it lies in $X^0$.
By Definition~I.\ref{I-def:eigendata}, $\ellst(0)(2\Re\phist(0))=\ellst(0)(\phist(0))+\ellst(0)(\overline{\phist(0)})=1$.
Hence
$|\ellst(0)(\Phi^\kappa-\Phi)-\kappa|\le\|\ellst(0)\|_{(X^0)^*}C'_0\kappa^2\le\frac\kappa2$, and
$|\ellst(0)(\Phi^\kappa-\Phi)|\in[\frac\kappa2,\frac{3\kappa}2]\subset[\frac\kappa2,2\kappa]$.
\end{proof}
